\documentclass[11pt]{article}
\usepackage{tcee}
\addbibresource{3A04.bib}
\tcetema[El pensamiento económico de Keynes]{3.A.4}{El pensamiento económico de Keynes. Referencia a la economía postkeynesiana y el desequilibrio.}
\tcefecha{04/11/2026}
\begin{document}
\tceetitulo

\section*{Introducción}
\minutos{2,5}

\begin{itemize}
\item \textbf{\emph{Enganche}}:
  \begin{itemize}
  \item «Pan para hoy, hambre para mañana», dice el refrán: lo que se gasta hoy sin pensar en el futuro se paga mañana. Es la objeción de siempre al gasto público financiado con deuda.
  \item En julio de 2020, con media economía parada por la pandemia, el Consejo Europeo acordó \emph{NextGenerationEU}: 750.000 millones de euros (a precios de 2018) financiados con deuda común para sostener la inversión y la demanda \parencite{CE2026NGEU}. Nadie propuso esperar a que bajaran los salarios.
  \item El propio nombre del fondo le da la vuelta al refrán: el pan de hoy se gasta pensando en la próxima generación, porque quien no come hoy no puede construir el futuro con el estómago vacío.
  \item Ese giro es de \autor{John Maynard Keynes}. A quienes confiaban en que la economía vuelve sola a su equilibrio les contestó: «este largo plazo es una mala guía para los asuntos de hoy. A largo plazo, todos muertos» \parencite[cap.~3, pág.~80; traducción propia]{Keynes1923Tract}.
  \item Detrás está la pregunta de esta exposición: \emph{¿se corrige sola una economía con paro, o hay que empujarla?}
  \end{itemize}

\item \textbf{\emph{Relevancia}}:
  \begin{itemize}
  \item \emph{Para la teoría}: con la \emph{Teoría general} nace la \concepto{macroeconomía} como disciplina, el estudio de la producción, el empleo y los precios de la economía en su conjunto, con los agregados (consumo, inversión, ahorro, renta) que mide la contabilidad nacional \vertema{3.A.30}.
  \item \emph{Para la política económica}: la política fiscal de estabilización y la idea de que el Estado responde del nivel de empleo vienen de \autor{Keynes}.
  \item \emph{Para la actualidad}: cada gran recesión lo ha devuelto al centro del debate. En 2008-2009 y en 2020 se discutió el tamaño de los multiplicadores fiscales \parencite{BlanchardLeigh2013Growth} \vertema{3.A.40}, qué hacer con los tipos de interés en cero, la vieja trampa de la liquidez \vertema{3.A.37}, y la inestabilidad financiera de \autor{Minsky} \vertema{3.A.42}.
  \item Además, como toda historia del pensamiento económico, enseña a no quedarse con el último manual y a ver por qué unas ideas llegan en unas épocas y no en otras: son 2 de los 4 motivos de \textcite{Schumpeter1954History}.\footnote{Los otros 2: la historia es fuente de ideas (principios que no fructificaron pueden orientar la investigación de hoy) y muestra cómo se desplazan las fronteras de la ciencia (cada época trae sus problemas y su forma de afrontarlos). La evolución del pensamiento económico resulta de la interacción de varios factores: el contexto económico de cada momento, el pensamiento anterior y los avances de otras disciplinas, como la filosofía, las matemáticas, la física o la biología. Conocerla permite entender las raíces del pensamiento actual, analizar los fenómenos económicos pasados y presentes y valorar los distintos programas de investigación.}
  \end{itemize}

\item \textbf{\emph{Contextualización}}:
  \begin{itemize}
  \item La historia del pensamiento económico avanza por \emph{revoluciones}: un \emph{paradigma} domina hasta que acumula anomalías que no sabe explicar y otro lo sustituye \parencite{Kuhn1962Structure} \vertema{3.A.1}.
  \item Para los \emph{clásicos} \vertema{3.A.2} y los \emph{neoclásicos} \vertema{3.A.3}, la oferta crea su propia demanda (la ley de \autor{Say}), los precios y los salarios flexibles vacían los mercados y el dinero solo fija el nivel de precios. \autor{Keynes} los llamará a todos «clásicos» (anexo~\ref{anexo:clasicos}).
  \item La anomalía fue el paro de entreguerras. En el Reino Unido, el paro de los trabajadores asegurados no bajó del 9,7~\% de 1921 a 1929 y llegó al 22,1~\% en 1932 \parencite{BoE2017Millennium}; en Estados Unidos, el PIB real cayó un 26~\% entre 1929 y 1933 \parencite{BEA2026NIPA} y el paro llegó al 25~\% de la población activa \parencite{Census1975HSUS}. La receta ortodoxa, bajar los salarios, no lo curaba.
  \item \textcite{Keynes1936} responde con la \emph{Teoría general de la ocupación, el interés y el dinero}: la producción y el empleo los fija la \concepto{demanda efectiva}, y una economía de mercado puede quedarse con \concepto{paro involuntario} sin que nada la saque de ahí.
  \item El libro es rico y oscuro a la vez, y sus herederos lo leyeron de 3 maneras:
    \begin{itemize}
    \item la \emph{síntesis neoclásica} lo tradujo al lenguaje de los neoclásicos con el modelo IS-LM: \autor{Keynes} para el corto plazo, con salarios rígidos, y los neoclásicos para el largo \vertema{3.A.5};
    \item los \emph{postkeynesianos} de Cambridge vieron en esa traducción una traición y llevaron la demanda efectiva y la incertidumbre al largo plazo, a la distribución y a las finanzas;
    \item los \emph{keynesianos del desequilibrio} releyeron a \autor{Keynes} como una teoría de los intercambios a precios que no vacían los mercados, y abrieron el camino a la nueva economía keynesiana \vertema{3.A.7}.
    \end{itemize}
  \end{itemize}

\item \textbf{\emph{Problemática}}:
  \begin{itemize}
  \item ¿Qué es la demanda efectiva y por qué puede haber paro involuntario aunque los salarios bajen?
  \item ¿Qué papel tienen el dinero, el tipo de interés y la incertidumbre?
  \item ¿Qué política económica se deduce?
  \item ¿Era la síntesis la única lectura posible de \autor{Keynes}? ¿Qué añaden Cambridge y el desequilibrio?
  \end{itemize}

\item \textbf{\emph{Estructura}}:
  \begin{itemize}
  \item La exposición va del modelo que \autor{Keynes} critica a su alternativa, y de esta a lo que hicieron con ella los herederos que no aceptaron la síntesis:
  \begin{itemize}
  \item \concepto{I. La escuela keynesiana y \autor{John Maynard Keynes}} (\ref{sec:keynes}): quiénes forman la escuela, el contexto y la obra de \autor{Keynes} y el modelo neoclásico del que parte;
  \item \concepto{II. El pensamiento económico de \autor{Keynes}: formalización y política económica} (\ref{sec:formalizacion}): la demanda efectiva, el dinero y el paro involuntario, lo que implican para la política y cómo se leyó la \emph{Teoría general};
  \item \concepto{III. La economía postkeynesiana y el desequilibrio} (\ref{sec:herederos}): la escuela de Cambridge (\autor{Kalecki}, \autor{Robinson}, \autor{Kaldor}, \autor{Minsky}) y los teóricos del desequilibrio (\autor{Patinkin}, \autor{Clower}, \autor{Leijonhufvud}, \autor{Malinvaud}).
  \end{itemize}
  \end{itemize}
\end{itemize}

\begin{esquema}[unbreakable]
\textbf{3.A.4. El pensamiento económico de \autor{Keynes}. Referencia a la economía postkeynesiana y el desequilibrio}\\[2pt]
Introducción\\
I. La escuela keynesiana y \autor{John Maynard Keynes}\\
\hspace*{1.5em}1. Caracterización general de la escuela keynesiana según \autor{Schumpeter}\\
\hspace*{1.5em}2. Contexto, método y obra\\
\hspace*{1.5em}3. El punto de partida: el modelo neoclásico\\
II. El pensamiento económico de \autor{Keynes}: formalización y política económica\\
\hspace*{1.5em}1. Formalización del pensamiento de \autor{Keynes} y comparación con el modelo neoclásico\\
\hspace*{1.5em}2. Implicaciones de política económica\\
\hspace*{1.5em}3. Valoración: las lecturas de la \emph{Teoría general}\\
III. La economía postkeynesiana y el desequilibrio\\
\hspace*{1.5em}1. La economía postkeynesiana: la escuela de Cambridge\\
\hspace*{1.5em}2. Los keynesianos del desequilibrio\\
Conclusión
\end{esquema}

\begin{notaopositor}
Reparto del cante: introducción, \tiempo{2,5}; I, \tiempo{6} (I.1, \tiempo{1,5}; I.2, \tiempo{2}; I.3, \tiempo{2,5}); II, \tiempo{10,5} (II.1, \tiempo{7,5}: bienes, \tiempo{3,5}, dinero, \tiempo{2}, y trabajo, \tiempo{2}; II.2, \tiempo{2}; II.3, \tiempo{1}); III, \tiempo{10} (III.1, \tiempo{5,5}; III.2, \tiempo{4,5}), y conclusión, \tiempo{1}. \autor{Keynes} se lleva \tiempo{16,5} y sus herederos \tiempo{10}: es la proporción que pide el título, que deja a los postkeynesianos y al desequilibrio como «referencia». La pizarra se organiza en 3 zonas. A la izquierda, el refrán del enganche y las 6 pegatinas, cada una junto a su concepto: «¿Dónde va Vicente? Donde va la gente» (las convenciones del capítulo~12), «puedes llevar el caballo al río, pero no obligarlo a beber» (la trampa de la liquidez), «desnudar a un santo para vestir a otro» (bajar los salarios, capítulo~19), «no se pueden sumar peras con manzanas» (la controversia del capital), «después de la calma viene la tempestad» (\autor{Minsky}) y «de donde no hay, no se puede sacar» (la decisión dual de \autor{Clower}). En el centro, los gráficos, que se borran al cambiar de bloque. A la derecha, el árbol de la escuela keynesiana, que se dibuja con la caracterización de la escuela y se completa con sus herederos, y, recuadradas y conservadas desde la formalización hasta el final, la demanda efectiva con su multiplicador y la preferencia por la liquidez. Lo que no se canta (el anexo~\ref{anexo:modelos} y los párrafos de profundización) sirve para las preguntas del tribunal.
\end{notaopositor}

% ==========================================================================
\section{La escuela keynesiana y \autor{John Maynard Keynes}}\label{sec:keynes}
\minutos{6}

\begin{itemize}
\item Empezamos por la escuela que fundó \autor{Keynes}, por el autor y por la teoría que quería sustituir.
\item \emph{Conclusión anticipada}: \autor{Keynes} funda una escuela cuyo objeto es explicar el paro involuntario como un fallo del sistema económico en su conjunto, y no de un mercado. Lo hace con el paro británico de entreguerras delante y frente a un modelo neoclásico en el que la flexibilidad de los precios, los salarios y el tipo de interés asegura el pleno empleo y el dinero es neutral.
\end{itemize}

% --------------------------------------------------------------------------
\subsection{Caracterización general de la escuela keynesiana según \autor{Schumpeter}}\label{sec:schumpeter}

\begin{itemize}
\item \autor{Schumpeter} reconoce una escuela de pensamiento por 3 rasgos: un \emph{líder}, unos \emph{seguidores} y un \emph{cuerpo doctrinal} propio, con su objeto y su método \parencite{Schumpeter1954History}. El keynesianismo los tiene:
  \begin{itemize}
  \item \emph{Líder}:
    \begin{itemize}
    \item \autor{John Maynard Keynes} (1883-1946).
    \end{itemize}
  \item \emph{Seguidores}:
    \begin{itemize}
    \item Se reparten en 3 ramas que discrepan entre sí sobre qué dijo de verdad \autor{Keynes}:
    \begin{itemize}
    \item la \emph{síntesis neoclásica}: \autor{John Hicks}, \autor{Alvin Hansen}, \autor{Paul Samuelson}, \autor{Franco Modigliani}, \autor{James Tobin}, \autor{Robert Solow}, \autor{Lawrence Klein} y \autor{James Meade} \vertema{3.A.5};
    \item los \emph{postkeynesianos} de Cambridge: \autor{Richard Kahn}, \autor{Joan Robinson}, \autor{Nicholas Kaldor}, \autor{Piero Sraffa} y \autor{Luigi Pasinetti}, con \autor{Michał Kalecki} como precursor y, en Estados Unidos, \autor{Hyman Minsky}, \autor{Paul Davidson} y \autor{Sidney Weintraub} (\ref{sec:postkeynesianos});
    \item los \emph{keynesianos del desequilibrio}: \autor{Don Patinkin} como antecedente, \autor{Robert Clower}, \autor{Axel Leijonhufvud}, \autor{Robert Barro}, \autor{Herschel Grossman} y \autor{Edmond Malinvaud} (\ref{sec:desequilibrio}).
    \end{itemize}
    \end{itemize}
  \item \emph{Cuerpo doctrinal}:
    \begin{itemize}
    \item \emph{objeto}: demostrar que puede existir \concepto{paro involuntario}, el de quien quiere trabajar al salario vigente y no encuentra empleo, y que su causa es una demanda agregada insuficiente: un fallo de la economía en su conjunto, no del mercado de trabajo \parencite{DeVroey2016History};\footnote{En el capítulo~2, \autor{Keynes} admite que la teoría clásica explica 2 clases de paro: el \emph{friccional}, por los desajustes temporales entre puestos y trabajadores, y el \emph{voluntario}, el de quien no acepta trabajar al salario que corresponde a su productividad marginal. Lo que no explica es una tercera clase, el paro \emph{involuntario} \parencite[cap.~2]{Keynes1936}.}
    \item \emph{método}: el análisis \emph{agregado} y de \emph{corto plazo} de un sistema de mercados interdependientes, con el instrumental de \autor{Marshall} (la competencia perfecta y el razonamiento verbal) y no con el equilibrio general de \autor{Walras};\footnote{En el Cambridge de los años 30 se apreciaba poco a \autor{Walras}, y \autor{Keynes} no veía útil su teoría para lo que quería demostrar \parencite{DeVroey2016History}. Se le atribuye un juicio muy duro sobre la teoría walrasiana en una carta de 1934 a \autor{Nicholas Georgescu-Roegen}. En el prefacio a la edición francesa, \textcite{Keynes1939Prefacio} explica que llama «general» a su teoría porque se ocupa sobre todo del comportamiento del sistema económico en su conjunto (la renta, el beneficio, la producción, el empleo y el ahorro agregados) y no del de una industria o un individuo.} un papel central de la \emph{incertidumbre} y de las expectativas; y la economía como ciencia moral, al servicio de la política.
    \end{itemize}
  \end{itemize}
\end{itemize}

\begin{table}[H]
\centering
\caption{Los supuestos de partida: los neoclásicos y \autor{Keynes}}\label{tab:supuestos}
\small
\begin{tabular}{>{\raggedright\arraybackslash}m{2.6cm}>{\raggedright\arraybackslash}m{5.2cm}>{\raggedright\arraybackslash}m{6cm}}
\toprule
\rowcolor{tceeazulpalido}\textbf{Supuesto} & \textbf{Neoclásicos} & \textbf{\autor{Keynes} (1936)}\\
\midrule
Horizonte & Largo plazo: el equilibrio al que tiende la economía & \textcolor{tceeazul}{\textbf{Corto plazo}}, con el equipo productivo dado: «a largo plazo, todos muertos» \parencite[cap.~3]{Keynes1923Tract}\\
\addlinespace[0.3em]
Información & Perfecta, o riesgo que se puede calcular & \textcolor{tceeazul}{\textbf{Incertidumbre}} que no se reduce a probabilidades: expectativas volátiles y convenciones (cap.~12)\\
\addlinespace[0.3em]
Salario & Real, fijado por la oferta y la demanda de trabajo; flexible & \textcolor{tceeazul}{\textbf{Nominal}}, negociado; los trabajadores defienden su salario relativo y se resisten a que se lo bajen (cap.~2); aunque bajara, no garantizaría el pleno empleo (cap.~19)\\
\addlinespace[0.3em]
Precios & Flexibles, en competencia perfecta & También flexibles, en competencia perfecta: suben con el empleo por los rendimientos decrecientes\\
\addlinespace[0.3em]
Tipo de interés & Real; iguala el ahorro y la inversión (fondos prestables) & \textcolor{tceeazul}{\textbf{Monetario}}; lo fija la preferencia por la liquidez (cap.~13)\\
\addlinespace[0.3em]
Dinero & Medio de cambio; neutral & \textcolor{tceeazul}{\textbf{También depósito de valor}}: puede atesorarse\\
\addlinespace[0.3em]
Ley de \autor{Say} & Se cumple: no hay exceso de oferta general & \textcolor{tceeazul}{\textbf{No se cumple}}: la demanda efectiva fija la producción (cap.~3)\\
\bottomrule
\end{tabular}
\fuente{elaboración propia a partir de \textcite{Keynes1936} y \textcite{DeVroey2016History}. En azul, los supuestos que cambia \autor{Keynes}.}
\end{table}

\begin{itemize}
\item \emph{Lo que cambia frente a los neoclásicos}: \autor{Keynes} mantiene la competencia perfecta y no recurre a la competencia imperfecta que se estaba formulando en esos años:\footnote{En 1933 aparecieron a la vez \emph{The Economics of Imperfect Competition}, de \autor{Joan Robinson} \parencite{Robinson1933Economics}, y \emph{The Theory of Monopolistic Competition}, de \autor{Edward Chamberlin} \parencite{Chamberlin1933Theory} \vertema{3.A.18}.} si el paro se debiera al poder de mercado o a los sindicatos, sería un problema de monopolios, y \autor{Keynes} buscaba un fallo más profundo, que se diera incluso con competencia perfecta \parencite{DeVroey2016History}. Cambia, en cambio, los supuestos de la tabla~\ref{tab:supuestos}, que el bloque~\ref{sec:formalizacion} desarrolla uno a uno.
\end{itemize}

% --------------------------------------------------------------------------
\subsection{Contexto, método y obra}\label{sec:contexto}

\begin{itemize}
\item \emph{Formación}:
  \begin{itemize}
  \item Nace en Cambridge en 1883, hijo del economista \autor{John Neville Keynes} y de \autor{Florence Ada Brown}.\footnote{\autor{Florence Ada Keynes} fue en 1932 la segunda mujer alcaldesa de Cambridge.}
  \item Estudia matemáticas en el King's College, donde se gradúa en 1905, y se forma como economista con \autor{Alfred Marshall} y \autor{Arthur Cecil Pigou}.
  \item De \autor{Marshall} toma el \emph{método}: el razonamiento de corto plazo, el equilibrio en tiempo real y el peso de las expectativas \vertema{3.A.3} \vertema{3.A.16}.
  \end{itemize}
\item \emph{Del Estado a Cambridge}:
  \begin{itemize}
  \item En 1906 entra en la India Office, de donde sale su primer libro, \emph{Indian Currency and Finance} \parencite{Keynes1913Indian}; en 1908 vuelve a Cambridge como profesor.\footnote{Al principio, el puesto lo pagaban de su bolsillo \autor{Marshall} y \autor{Pigou}. De 1911 a 1945 dirigió \emph{The Economic Journal}; en 1925 se casó con la bailarina rusa \autor{Lydia Lopokova}, y formó parte del grupo de Bloomsbury, con \autor{Virginia Woolf} y \autor{Lytton Strachey}. En 1942 recibió el título de barón \autor{Keynes} de Tilton.}
  \item Su disertación sobre la probabilidad le da en 1909 la \emph{fellowship} del King's College y se publica en 1921 como \emph{A Treatise on Probability} \parencite{Keynes1921Treatise}.
  \item La probabilidad como grado de creencia razonable, que no siempre puede medirse, es la raíz de su \emph{incertidumbre} \vertema{3.A.10}.
  \end{itemize}
\item \emph{Las 3 etapas de su obra}:
  \begin{itemize}
  \item Cada etapa está ligada a un problema de su tiempo \parencite[cap.~14]{Roncaglia2006Riqueza} (gráfico~\ref{fig:cronologia}):
  \begin{itemize}
  \item \emph{la Primera Guerra Mundial y sus secuelas}: trabaja en el Tesoro durante la guerra;\footnote{En 1917 recibió la Orden del Baño, en el grado de \emph{Companion} (CB), no el de caballero.} representa al Tesoro en la Conferencia de Paz de París, dimite y critica las reparaciones impuestas a Alemania en \emph{The Economic Consequences of the Peace} \parencite{Keynes1919Economic}; en el \emph{Tract on Monetary Reform} \parencite{Keynes1923Tract} defiende la estabilidad de los precios frente a la vuelta al patrón oro a la paridad de antes de la guerra, y en \emph{The End of Laissez-Faire} \parencite{Keynes1926End} anuncia que el Estado tendrá que hacer lo que nadie hace;
  \item \emph{la Gran Depresión}: el \emph{Treatise on Money} \parencite{Keynes1930Treatise}, todavía dentro de la teoría cuantitativa, se discute en 1930-1931 en el «círculo» de jóvenes economistas de Cambridge (\autor{Richard Kahn}, \autor{Joan Robinson} y \autor{Austin Robinson}, \autor{Piero Sraffa} y \autor{James Meade}), la semilla de la escuela postkeynesiana (\ref{sec:postkeynesianos}). De ese círculo sale el multiplicador de \textcite{Kahn1931Relation}, y de la discusión, la \emph{Teoría general} \parencite{Keynes1936}, que \autor{Keynes} defiende al año siguiente frente a sus críticos \parencite{Keynes1937General};
  \item \emph{la Segunda Guerra Mundial}: en \emph{How to Pay for the War} \parencite{Keynes1940How} propone el ahorro obligatorio para financiar la guerra sin inflación. En Bretton Woods (1944) presenta el \emph{Plan \autor{Keynes}}, una unión de compensación con una moneda internacional, el \emph{bancor}, y un ajuste simétrico entre acreedores y deudores; se impone el \emph{Plan \autor{White}} estadounidense \vertema{3.B.20}. En marzo de 1946 asiste en Savannah, como gobernador por el Reino Unido, a la reunión inaugural del FMI y del Banco Mundial \parencite{Hansard1946Savannah} \vertema{3.B.22}, y muere el 21 de abril en Tilton (Sussex).
  \end{itemize}
  \end{itemize}
\end{itemize}

\begin{figure}[H]
\centering
\caption{La vida de \autor{John Maynard Keynes} (1883-1946)}\label{fig:cronologia}
\grafico[0.95\textwidth]{cronologia-keynes}
\fuente{elaboración propia a partir de \textcite{Skidelsky1983Hopes,Skidelsky1992Saviour,Skidelsky2000Fighting,Skidelsky2003Keynes}. Las filas no están a escala. A la izquierda, los 3 volúmenes de la biografía de \autor{Robert Skidelsky}.}
\end{figure}

\begin{anotaciones}
\textbf{\autor{Keynes}, según las biografías de \autor{Robert Skidelsky}}
\begin{itemize}
\item \emph{La biografía}: 3 volúmenes \parencite{Skidelsky1983Hopes,Skidelsky1992Saviour,Skidelsky2000Fighting} y una versión abreviada en uno \parencite{Skidelsky2003Keynes}. Los títulos resumen la vida: las esperanzas traicionadas en Versalles (1883-1920), el economista como salvador (1920-1937) y la lucha por Gran Bretaña durante la guerra (1937-1946).
\item \emph{Eton y Cambridge}: entra en Eton con beca en 1897 y llega al King's College con otra. En 1903 lo eligen para los Apóstoles, la sociedad secreta de debate de Cambridge, donde le marcó la ética de \autor{George Edward Moore} \parencite{Skidelsky1983Hopes}.
\item \emph{El examen de 1906}: queda segundo en el examen de ingreso en la administración y su peor nota es la de economía; comentó que, por lo visto, los examinadores sabían menos que él. El primero eligió el Tesoro, y \autor{Keynes} tuvo que conformarse con la India Office \parencite{Skidelsky1983Hopes}.
\item \emph{Las pesetas de 1915}: ya en el Tesoro, se cuenta que reunió una gran cantidad de pesetas y las vendió de golpe para bajar la cotización de la moneda española, lo que convenía al Reino Unido, que compraba mucho a una España neutral \parencite[pág.~34]{Lekachman1966Age}.
\item \emph{El especulador}: en 1920 especuló con divisas y estuvo cerca de la quiebra; salió adelante con préstamos de sus amigos y llegó a amasar una fortuna. Como tesorero (\emph{bursar}) del King's College gestionó durante décadas sus inversiones \parencites{Skidelsky1992Saviour}{Accominotti2016Smart}.
\item \emph{Las artes}: en febrero de 1936, el mismo mes que la \emph{Teoría general}, abre el Arts Theatre de Cambridge, pagado en gran parte de su bolsillo; durante la guerra preside el comité de fomento de la música y las artes, y lo eligen para presidir el Arts Council de Gran Bretaña, que nace en 1946, el año de su muerte \parencite{Skidelsky2000Fighting}.
\item \emph{\autor{Newton}}: en 1936 compra en la subasta de Sotheby's manuscritos de \autor{Isaac Newton}, muchos de alquimia, y los reúne poco a poco. En su conferencia «Newton, the Man», que leyó su hermano \autor{Geoffrey Keynes} en 1946, tras su muerte, lo describe como el último de los magos y no el primero de la edad de la razón \parencite{Skidelsky2000Fighting}.
\item \emph{La enfermedad}: en 1937 sufre un infarto; lo trata el médico húngaro \autor{Janos Plesch}, y desde entonces trabaja con el corazón enfermo hasta su muerte en 1946, agotado por las negociaciones de Bretton Woods y del préstamo estadounidense de 1945 \parencite{Skidelsky2000Fighting}.
\end{itemize}
\end{anotaciones}

\begin{itemize}
\item \emph{El contexto económico}:
  \begin{itemize}
  \item \autor{Keynes} escribe con el paro británico delante, no solo con el \emph{crash} de 1929 (gráfico~\ref{fig:reino-unido}).
  \item En 1921 el PIB real británico cayó un 9,7~\% y el paro de los asegurados llegó al 16,9~\%; en toda la década no volvió a bajar del 9,7~\% \parencite{BoE2017Millennium}.
  \item La vuelta al patrón oro a la paridad de antes de la guerra, en abril de 1925, encareció la libra y obligaba a bajar los salarios y los precios para recuperar la competitividad.
  \item Con la Gran Depresión, el Reino Unido abandonó el oro en septiembre de 1931, y el paro de los asegurados subió al 22,1~\% en 1932 \parencite{BoE2017Millennium}.
  \item Para la teoría de la época el paro tenía que ser pasajero: si persistía, era por algo que impedía bajar los salarios (\ref{sec:neoclasicos}).
  \end{itemize}
\end{itemize}

\begin{figure}[H]
\centering
\caption{Reino Unido, 1920-1938: paro de los trabajadores asegurados y PIB real}\label{fig:reino-unido}
\grafico[0.85\textwidth]{reino-unido-entreguerras}
\fuente{elaboración propia a partir de \textcite{BoE2017Millennium}. El paro es la tasa de los trabajadores asegurados, la cifra oficial de la época (unos 12 millones de trabajadores, sobre todo manuales); sobre la población activa total, comparable con la EPA, fue del 15,4~\% en 1932. PIB real en índice 1929 = 100.}
\end{figure}

\begin{itemize}
\item \emph{Su propósito}:
  \begin{itemize}
  \item No era un revolucionario: quería salvar el capitalismo de sí mismo, en una época en la que el paro masivo y el ascenso de los totalitarismos lo ponían en peligro.
  \item Proponía un capitalismo reformado, con el Estado responsable del nivel de empleo, que conservara la libertad y la eficiencia de los mercados \parencite{Skidelsky2003Keynes}.
  \end{itemize}
\end{itemize}

% --------------------------------------------------------------------------
\subsection{El punto de partida: el modelo neoclásico}\label{sec:neoclasicos}

\paragraph{a) Idea}

\begin{itemize}
\item Para entender qué cambia \autor{Keynes}, hay que ver antes el modelo que quería sustituir: el punto de partida (\emph{baseline}) frente al que se mide la \emph{Teoría general}.
\item Ese modelo no lo escribió un solo autor. Es la reconstrucción habitual a partir de \autor{Marshall}, \autor{Pigou}, \autor{Fisher} y \autor{Wicksell} \parencite[cap.~1]{Heijdra2017Foundations}.
\item De \autor{Marshall} viene el mercado que se vacía: en su mercado de pescado, el precio vacía el mercado cada día y, con tiempo, la cantidad se ajusta hacia el equilibrio normal \parencite[libro~V]{Marshall1890Principles} \vertema{3.A.16} (anexo~\ref{anexo:modelos}). En un mercado así no hay recursos ociosos persistentes. \autor{Keynes} conservará el método de su maestro, el corto plazo y el ajuste por cantidades, pero lo llevará a la economía entera.
\end{itemize}

\begin{anotaciones}
El «modelo neoclásico» de este apartado y el «modelo keynesiano» del bloque~\ref{sec:formalizacion} son reconstrucciones: reúnen aportaciones de distintos autores y usan gráficos y fórmulas que \autor{Keynes} y sus maestros no siempre escribieron. Además, el pensamiento de \autor{Keynes} admite varias lecturas; lo que sigue es la formalización habitual, con la notación del resto del temario.
\end{anotaciones}

\paragraph{b) Modelo}

\subparagraph{i) Supuestos}

\begin{itemize}
\item Economía cerrada, con el capital dado, en la que los mercados de bienes, de trabajo y de dinero son de competencia perfecta.
\item Precios, salarios y tipo de interés \emph{flexibles}: cada mercado se vacía.
\item Información perfecta y agentes racionales que deciden en términos \emph{reales}, sin ilusión monetaria.
\item El dinero solo es medio de cambio: nadie guarda dinero ocioso pudiendo prestarlo con interés.
\end{itemize}

\subparagraph{ii) Desarrollo}

\begin{itemize}
\item La causalidad va del mercado de trabajo al de dinero: el trabajo fija el empleo y la producción, los fondos prestables reparten esa producción entre consumo e inversión, y el dinero fija los precios.
\item \emph{El mercado de trabajo fija el empleo} \vertema{3.A.25}. Las empresas contratan hasta que el salario real iguala la productividad marginal del trabajo, y los hogares ofrecen más trabajo cuanto mayor es el salario real:
  \begin{equation}
  w\equiv\frac{W}{P}=\mathrm{PMg}_L(L),\qquad L^S=L^S(w),\qquad L^D(w^c)=L^S(w^c)\ \Rightarrow\ L^c,\quad Y_F=F(L^c). \label{eq:trabajo-neoclasico}
  \end{equation}
  \begin{itemize}
  \item Con el salario real flexible, el mercado se vacía en $(L^c,w^c)$, el pleno empleo, y la función de producción da la producción de pleno empleo, $Y_F$ (panel a del gráfico~\ref{fig:modelo-neoclasico}).
  \item La \emph{oferta agregada es vertical} en $Y_F$: si suben los precios, el salario nominal sube en la misma proporción, el real no cambia y tampoco el empleo.
  \end{itemize}
\item \emph{El mercado de bienes fija el tipo de interés real} \vertema{3.A.33}. Con la producción dada, la igualdad entre la producción y el gasto, $Y=C+I+G$, se lee como un mercado de \concepto{fondos prestables} (gráfico~\ref{fig:fondos-prestables}):
  \begin{equation}
  S(r)=I(r)+G-T\ \Rightarrow\ r^n. \label{eq:fondos-prestables}
  \end{equation}
  \begin{itemize}
  \item La \emph{oferta} de fondos es el ahorro, que crece con el tipo de interés real (premia la renuncia al consumo presente); la \emph{demanda}, la inversión más el déficit público, que decrece con él (se invierte hasta que la productividad marginal del capital iguala el tipo).
  \item El tipo que iguala las 2 es el \emph{tipo natural}, $r^n$, de \autor{Wicksell}.\footnote{La idea de un mercado de fondos que se prestan y se toman prestados es de \autor{Marshall} y \autor{Wicksell}; la expresión «fondos prestables» es de \autor{Robertson} y \autor{Ohlin}, en los años 30, ya en polémica con \autor{Keynes} \parencite{DeVroey2016History}. Para \autor{Wicksell}, el tipo natural puede diferir del que cobran los bancos, y esa diferencia mueve los precios.}
  \end{itemize}
\end{itemize}

\begin{figure}[H]
\centering
\caption{El mercado de fondos prestables neoclásico}\label{fig:fondos-prestables}
\grafico[0.55\textwidth]{fondos-prestables}
\fuente{elaboración propia a partir de \textcite{Fonseca2026Neoclassical}.}
\end{figure}

\begin{itemize}
\item \emph{El mercado de dinero fija el nivel de precios}: la \concepto{teoría cuantitativa} \vertema{3.A.35}. Con la ecuación de cambio de \textcite{Fisher1911Purchasing}, $M\,V=P\,Y$, o con la demanda de saldos de caja de Cambridge, $M^D=k\,P\,Y$ \parencite{Pigou1917Value}, con $k=1/V$, y con $V$ (o $k$) estable y la producción en $Y_F$:
  \begin{equation}
  P=\frac{M\,V}{Y_F}=\frac{M}{k\,Y_F}. \label{eq:cuantitativa-clasica}
  \end{equation}
  \begin{itemize}
  \item Para cada cantidad de dinero, $P\,Y=M\,V$ es una \emph{demanda agregada} con forma de hipérbola en el plano $(Y,P)$. Con la oferta agregada vertical, más dinero solo sube los precios en la misma proporción (panel b del gráfico~\ref{fig:modelo-neoclasico}).
  \end{itemize}
\end{itemize}

\begin{figure}[H]
\centering
\caption{El modelo neoclásico: el mercado de trabajo fija el empleo y la producción; el dinero, el nivel de precios}\label{fig:modelo-neoclasico}
\grafico[0.9\textwidth]{modelo-neoclasico}
\fuente{elaboración propia a partir de \textcite{Fonseca2026Neoclassical} y \textcite[cap.~1]{Heijdra2017Foundations}; la misma escala que el mercado de trabajo keynesiano (gráfico~\ref{fig:keynesiano-shock}).}
\end{figure}

\subparagraph{iii) Implicaciones}

\begin{itemize}
\item \emph{Se cumple la ley de \autor{Say}}: todo lo que se ahorra se invierte, porque el tipo de interés real iguala el ahorro y la inversión. No puede haber un exceso de oferta general de bienes.
\item \emph{La política fiscal no mueve la producción}: más gasto público sube el tipo de interés y desplaza la misma cantidad de gasto privado (menos consumo, porque se ahorra más, y menos inversión). La \concepto{expulsión} es total: cambia la composición de la demanda, no su nivel. De ahí la defensa del presupuesto equilibrado.
\item \emph{La política monetaria tampoco}: el dinero es \concepto{neutral} y se cumple la \concepto{dicotomía clásica}: las variables reales se determinan sin el dinero.
\item \emph{El paro}: solo hay paro \emph{friccional} y \emph{voluntario}. Si persiste, es porque algo impide que el salario baje:
  \begin{itemize}
  \item los sindicatos, como estudia \textcite{Hicks1932Theory} \vertema{3.A.27};
  \item el poder de mercado de las grandes empresas, el de la competencia monopolística de \textcite{Chamberlin1933Theory} \vertema{3.A.18}.
  \end{itemize}
  La receta es flexibilizar: bajar los salarios y dejar hacer al mercado (\emph{laissez-faire}).
\end{itemize}

\paragraph{c) Valoración}

\begin{itemize}
\item Es un modelo coherente, en el que cada mercado fija una variable y la economía vuelve sola al pleno empleo.
\item Pero no explicaba un paro masivo que duraba años ni por qué bajar los salarios no lo curaba. Esa es la anomalía de la que parte \autor{Keynes}.
\item \emph{Recapitulación}: en el modelo neoclásico manda la oferta: el trabajo fija la producción, el tipo de interés real reparte el gasto y el dinero solo fija los precios. Una vez visto el modelo que \autor{Keynes} quería sustituir, pasamos a su alternativa. El tipo de interés será la primera víctima: de $r$ pasamos a $i$, y el interés deja de ser el precio del ahorro para ser el precio de la liquidez.
\end{itemize}

% ==========================================================================
\section{El pensamiento económico de \autor{Keynes}: formalización y\newline política económica}\label{sec:formalizacion}
\minutos{10,5}

\begin{itemize}
\item Una vez visto el modelo de partida, pasamos a la formalización de la \emph{Teoría general}, a sus consecuencias para la política económica y a las lecturas que se hicieron de ella.
\item \emph{Conclusión anticipada}: \autor{Keynes} invierte la causalidad neoclásica. La demanda efectiva fija la producción y el empleo; el tipo de interés se fija en el mercado de dinero, no en el de fondos prestables; el ahorro se ajusta a la inversión a través de la renta, y puede haber paro involuntario aunque los salarios bajen. El remedio es la política fiscal.
\end{itemize}

\begin{anotaciones}
\autor{Keynes} no escribió el modelo que sigue: la \emph{Teoría general} tiene pocas ecuaciones, y ya en el prefacio dirigía el libro a sus colegas economistas, a los que dejaba la tarea de juzgarlo \parencite{Keynes1936}. Las fórmulas son la forma lineal que se hizo habitual después, con la notación del resto del temario; donde la forma lineal no es fiel al libro, se avisa.
\end{anotaciones}

% --------------------------------------------------------------------------
\subsection{Formalización del pensamiento de \autor{Keynes} y comparación con el modelo neoclásico}\label{sec:formalizacion-keynes}

\begin{itemize}
\item \autor{Keynes} \emph{invierte la causalidad}: donde los neoclásicos empezaban por el mercado de trabajo, él termina. Por eso los mercados van en el orden de la \emph{Teoría general}:
  \begin{itemize}
  \item la demanda efectiva fija la producción (bienes);
  \item el dinero fija el tipo de interés, del que depende la inversión (dinero);
  \item el empleo sale de los 2, y puede quedarse por debajo del pleno empleo (trabajo).
  \end{itemize}
\end{itemize}

\subsubsection{El mercado de bienes: la demanda efectiva}\label{sec:bienes}

\paragraph{a) Idea}

\begin{itemize}
\item El capítulo~3 es el corazón del libro. Los empresarios dan el empleo con el que esperan maximizar su beneficio, y para eso comparan 2 funciones del empleo, $L$ \parencite[cap.~3]{Keynes1936}:\footnote{\autor{Keynes} escribe $N$ para el empleo: $Z=\phi(N)$ y $D=f(N)$. Aquí, $L$, como en el resto del temario.}
  \begin{itemize}
  \item la \emph{función de oferta agregada}, $Z(L)$: lo que tienen que esperar ingresar para que les compense dar ese empleo;
  \item la \emph{función de demanda agregada}, $D(L)$: lo que esperan ingresar vendiendo lo que produce ese empleo.
  \end{itemize}
\item El empleo se fija donde se cortan, $D(L)=Z(L)$: es el \concepto{punto de demanda efectiva}. Nada asegura que sea el de pleno empleo. La ley de \autor{Say} equivale a suponer que $D(L)=Z(L)$ para cualquier $L$: entonces todo empleo sería un equilibrio y la competencia lo llevaría hasta el máximo. Si no es así, la economía puede pararse antes.
\item La clave está en el consumo: cuando el empleo y la renta suben, el consumo sube menos, y la diferencia la tiene que cubrir la inversión. Si no la cubre, los empresarios no venden lo que esperaban y reducen el empleo.
\end{itemize}

\paragraph{b) Modelo}

\subparagraph{i) Supuestos}

\begin{itemize}
\item Economía cerrada a corto plazo, con el equipo productivo, la técnica y la población dados.
\item Competencia perfecta en el mercado de bienes y el salario nominal dado en cada momento.
\item Hay recursos ociosos: las empresas producen lo que se les demanda.
\item El tipo de interés no se fija en este mercado: viene del de dinero (\ref{sec:dinero}). Aquí se toma como dado, $i=i_0$.
\end{itemize}

\subparagraph{ii) Desarrollo}

\begin{itemize}
\item \emph{El consumo} \vertema{3.A.33}: la propensión a consumir es «una función bastante estable» de la renta \parencite[cap.~8; traducción propia]{Keynes1936}, según una \concepto{ley psicológica fundamental}: cuando la renta sube, el consumo sube, pero menos \parencite[cap.~3]{Keynes1936}. En forma lineal,
  \begin{equation}
  C_t=\overline{C}+c\,Y^d_t,\qquad 0<c<1,\qquad Y^d_t=(1-t)\,Y_t, \label{eq:consumo}
  \end{equation}
  con $\overline{C}$ el consumo autónomo, $c$ la propensión marginal a consumir y $t$ el tipo impositivo.
  \begin{itemize}
  \item La propensión media, $C_t/Y^d_t=\overline{C}/Y^d_t+c$, es mayor que la marginal y baja con la renta: quien gana más ahorra una parte mayor.
  \item El ahorro es lo que queda, $S_t=Y^d_t-C_t$, y ya no depende del tipo de interés, sino de la renta.\footnote{\autor{Keynes} escribe el consumo como función del empleo, $C=\chi(N)$, y define el ahorro como el exceso de la renta sobre el consumo \parencite[caps.~3 y 6]{Keynes1936}.}
  \end{itemize}
\end{itemize}

\begin{anotaciones}
\textbf{Una advertencia}: \autor{Keynes} no supuso que la propensión marginal a consumir fuera constante: «la propensión marginal a consumir no es constante para todos los niveles de empleo, y es probable que, por regla general, tienda a disminuir a medida que el empleo aumenta» \parencite[cap.~10, apdo.~III; traducción propia]{Keynes1936}. La recta \eqref{eq:consumo}, con $c$ constante, es una simplificación posterior de la síntesis neoclásica.
\end{anotaciones}

\begin{itemize}
\item \emph{La inversión} \vertema{3.A.34}: no depende del ahorro disponible, sino de lo que se espera ganar con ella y del tipo de interés \parencite[caps.~11 y 12]{Keynes1936}.
  \begin{itemize}
  \item La \concepto{eficiencia marginal del capital} (EMgK) es el tipo de descuento que iguala el valor actual de los rendimientos esperados de un bien de capital, $R^e_t$, con su precio de oferta, $P^K$:
    \begin{equation}
    P^K=\sum_{t=1}^{T}\frac{R^e_t}{(1+\mathrm{EMgK})^t}. \label{eq:emgk}
    \end{equation}
  \item Se invierte hasta que la EMgK de la última inversión iguala el tipo de interés, $\mathrm{EMgK}=i$. En forma lineal,
    \begin{equation}
    I=\overline{I}-b\,i,\qquad b>0, \label{eq:inversion}
    \end{equation}
    con $\overline{I}$ lo que depende de las expectativas y $b$ la sensibilidad al tipo de interés, que \autor{Keynes} creía pequeña.
  \item \emph{La incertidumbre} (capítulo~12): los rendimientos esperados dependen de un futuro que no se puede calcular. Como no hay base para hacerlo, los inversores siguen \emph{convenciones}: suponen que lo de hoy seguirá y, en bolsa, tratan de adivinar lo que pensará la mayoría, como en un concurso de belleza en el que se premia acertar la cara que elegirán los demás \parencite[cap.~12, apdo.~V]{Keynes1936}. Pegatina: «¿Dónde va Vicente? Donde va la gente».
  \item Cuando la convención se rompe, las expectativas se hunden. Por eso la inversión es la parte más \emph{inestable} de la demanda, movida por los \concepto{\emph{animal spirits}}: «un impulso espontáneo a la acción, en vez de a la inacción, y no el resultado de una media ponderada de beneficios cuantitativos multiplicados por probabilidades cuantitativas» \parencite[cap.~12, apdo.~VII; traducción propia]{Keynes1936}. Al responder a sus críticos, \textcite{Keynes1937General} dejó claro que la incertidumbre era el mensaje central del libro \vertema{3.A.10}.
  \end{itemize}
\item \emph{El equilibrio y el multiplicador}: la demanda agregada planeada es $\mathrm{DA}=\overline{C}+c\,(1-t)\,Y+\overline{I}-b\,i_0+\overline{G}$, con $\overline{G}$ el gasto público, y la producción se ajusta hasta igualarla, $Y=\mathrm{DA}$. A tipo de interés dado,
  \funcionclave[eq:multiplicador]{Y^*=\frac{\overline{C}+\overline{I}-b\,i_0+\overline{G}}{1-c\,(1-t)}\qquad\Longrightarrow\qquad \Delta Y^*=\frac{\Delta\overline{G}}{1-c\,(1-t)}}
  \begin{itemize}
  \item Es el \concepto{multiplicador}: cada euro de gasto autónomo es renta de alguien, que gasta una parte, que es renta de otro, y así sucesivamente. Como $0<c\,(1-t)<1$, el multiplicador es mayor que 1.
  \item \textcite{Kahn1931Relation} lo formuló para el empleo que crean las obras públicas, y \autor{Keynes} lo aplicó a la inversión \parencite[cap.~10]{Keynes1936}.
  \item Si el tipo de interés cambia con la renta, la misma ecuación da la curva IS de la síntesis \vertema{3.A.5}.
  \end{itemize}
\end{itemize}

\begin{figure}[H]
\centering
\caption{La cruz keynesiana: la producción se ajusta a la demanda planeada}\label{fig:cruz}
\grafico[0.55\textwidth]{cruz-keynesiana}
\fuente{elaboración propia a partir de \textcite{Samuelson1948Economics} y \textcite[caps.~3 y 10]{Keynes1936}.}
\end{figure}

\begin{itemize}
\item \emph{El ajuste} (la \concepto{cruz keynesiana}, en la forma que le dio \textcite{Samuelson1948Economics}; gráfico~\ref{fig:cruz}):
  \begin{itemize}
  \item a la derecha de $Y^*$ se produce más de lo que se demanda: lo que no se vende se acumula en existencias, una inversión que nadie había planeado (la inversión \emph{realizada supera} a la planeada), y las empresas reducen la producción;
  \item a la izquierda, las existencias bajan (la inversión realizada es \emph{menor} que la planeada) y las empresas producen más.
  \end{itemize}
  La variable que se ajusta es la \emph{renta}, no el tipo de interés: el ahorro se iguala a la inversión porque la renta se mueve hasta que el ahorro de los hogares es justo el que absorbe la inversión.
\end{itemize}

\subparagraph{iii) Implicaciones}

\begin{itemize}
\item \emph{La ley de \autor{Say} no se cumple}: nada asegura que $Y^*$ sea la producción de pleno empleo, $Y_F$. Si la inversión cae, la renta cae varias veces más.
\item \emph{La \concepto{paradoja del ahorro}}: si todos intentan ahorrar más, el consumo cae, la renta cae con el multiplicador y, en el nuevo equilibrio, el ahorro nacional, $S+(T-\overline{G})$, sigue siendo igual a la inversión, $\overline{I}-b\,i_0$: se ahorra lo mismo en total, con menos renta y menos empleo. Si la inversión depende de la renta, el ahorro total incluso baja. Lo que es prudente para uno es dañino para todos: \autor{Keynes} recuerda la \emph{Fábula de las abejas} de \autor{Bernard Mandeville} \parencite[cap.~23]{Keynes1936}.
\item \emph{Frente a los neoclásicos}: allí más ahorro bajaba el tipo de interés y subía la inversión; aquí la inversión apenas responde al tipo ($b$ pequeño) y ese mecanismo no compensa la caída de la demanda. El ahorro se ajusta a la inversión, y no al revés.
\end{itemize}

\subsubsection{El mercado de dinero: la preferencia por la liquidez}\label{sec:dinero}

\paragraph{a) Idea}

\begin{itemize}
\item Con la inversión ya no atada al ahorro, falta saber qué fija el tipo de interés. Para \autor{Keynes} no es el precio que iguala el ahorro y la inversión: es el precio de desprenderse de la liquidez, la recompensa por renunciar a ella y no por renunciar al consumo \parencite[cap.~13]{Keynes1936}.
\end{itemize}

\paragraph{b) Modelo}

\subparagraph{i) Supuestos}

\begin{itemize}
\item Dos activos financieros: el dinero, que no rinde, y los bonos a largo plazo, cuyo precio baja cuando sube el tipo de interés.
\item La oferta de dinero, $M$, es \emph{exógena}: la fija el banco central.
\item Cada inversor tiene su propia idea del tipo de interés «normal» al que volverá el mercado.
\end{itemize}

\subparagraph{ii) Desarrollo}

\begin{itemize}
\item \emph{La demanda de dinero} \vertema{3.A.35}: el público guarda dinero por 3 motivos \parencite[caps.~13 y 15]{Keynes1936}:
  \begin{itemize}
  \item \emph{transacción}: para los pagos corrientes, que crecen con la renta;
  \item \emph{precaución}: para imprevistos, también ligado a la renta;
  \item \emph{especulación}: si el tipo de mercado está por debajo del que un inversor cree normal, espera que suba y que baje el precio de los bonos, y prefiere el dinero para no perder; si está por encima, espera que baje y compra bonos para ganar con la subida de su precio. Como las ideas de lo «normal» son distintas, la demanda agregada de dinero especulativo sube de forma continua cuando baja el tipo.
  \end{itemize}
  Los 2 primeros motivos dependen de la renta, $L_1(Y)$, y el tercero, del tipo de interés, $L_2(i)$. En forma lineal:
  \funcionclave[eq:liquidez]{\frac{M}{P}=L_1(Y)+L_2(i)=k\,Y-h\,i,\qquad k,\,h>0}
  con $k$ la sensibilidad a la renta y $h$ al tipo de interés.\footnote{Si $h=0$, la demanda de dinero no depende del tipo y vuelve la de Cambridge, $M/P=k\,Y$: es el caso clásico. Si $h\to\infty$, la demanda es horizontal: es la trampa de la liquidez \vertema{3.A.5}.}
\item \emph{El tipo de interés} es el que iguala la demanda de saldos reales con la oferta (gráfico~\ref{fig:liquidez}):
  \begin{equation}
  i=\frac{1}{h}\Big(k\,Y-\frac{M}{P}\Big). \label{eq:tipo-keynes}
  \end{equation}
  Es un fenómeno \emph{monetario}: más dinero lo baja y más renta lo sube.
\item \emph{La \concepto{trampa de la liquidez}}: con un tipo tan bajo que todos esperan que suba, cualquier dinero de más se atesora y el tipo no baja: la demanda de dinero se vuelve horizontal \parencite[cap.~15, apdo.~II]{Keynes1936}. \autor{Keynes} la veía posible, aunque decía no conocer ningún caso. Pegatina: «puedes llevar el caballo al río, pero no obligarlo a beber».
\end{itemize}

\begin{figure}[H]
\centering
\caption{La preferencia por la liquidez: del inversor individual a la demanda agregada de dinero y la trampa de la liquidez}\label{fig:liquidez}
\grafico[0.85\textwidth]{preferencia-liquidez}
\fuente{elaboración propia a partir de \textcite[caps.~13 y 15]{Keynes1936}. La forma lineal $k\,Y-h\,i$ no tiene tramo horizontal: la trampa es el límite $h\to\infty$.}
\end{figure}

\subparagraph{iii) Implicaciones}

\begin{itemize}
\item \emph{El ataque a la ley de \autor{Say}}: si el dinero es también depósito de valor, guardarlo no es irracional. Una parte de la renta puede quedarse en saldos ociosos en lugar de gastarse, y la economía se atasca por debajo del pleno empleo:
  \begin{itemize}
  \item la ley de \autor{Say}, en la formulación de \textcite{Lange1942Say}, afirma que el valor de los excesos de demanda de los bienes, sin contar el dinero, es nulo: nunca puede haber un exceso de oferta general de bienes;
  \item en una economía monetaria, en cambio, un exceso de oferta de bienes puede ir acompañado de un exceso de demanda de dinero \parencite{BeckerBaumol1952Classical} \vertema{3.A.21}.
  \end{itemize}
  Es la ruptura más profunda con los neoclásicos, y la que acerca a \autor{Keynes} a los mercantilistas, a los que reivindica en el capítulo~23 por su preocupación por la escasez de dinero y por el tipo de interés \parencite{Barber1967History}.
\item \emph{El dinero no es neutral} a corto plazo y se rompe la dicotomía clásica: más $M$ baja el tipo, sube la inversión y, con el multiplicador, la renta.
\item Pero la cadena tiene 2 eslabones débiles: la trampa de la liquidez ($h$ muy grande: el tipo apenas baja) y una inversión que responde poco al tipo ($b$ pequeño: lo poco que baja mueve poco la inversión).
\item La demanda de dinero es \emph{inestable}, porque depende de expectativas sobre el tipo de interés que cambian con facilidad.
\item \emph{Frente a los neoclásicos}: el tipo de interés deja de ser real y de igualar el ahorro y la inversión; el dinero deja de ser solo medio de cambio.
\end{itemize}

\subsubsection{El mercado de trabajo: el paro involuntario}\label{sec:trabajo}

\paragraph{a) Idea}

\begin{itemize}
\item El empleo llega al final de la cadena: la demanda efectiva fija la producción, y la producción, el empleo. El mercado de trabajo no decide; recibe lo que deciden los otros 2.
\end{itemize}

\paragraph{b) Modelo}

\subparagraph{i) Supuestos}

\begin{itemize}
\item \emph{Lo que acepta}: la demanda de trabajo neoclásica, el «primer postulado clásico»: las empresas contratan hasta que el salario real iguala la productividad marginal del trabajo \parencite[cap.~2]{Keynes1936}.
\item \emph{Lo que rechaza}: el «segundo postulado», que el salario real iguale la desutilidad marginal del trabajo, es decir, que los trabajadores estén siempre sobre su curva de oferta.
\item \emph{El salario nominal}: los trabajadores negocian el salario nominal, $W$, y no pueden fijar el real, $W/P$, que depende de los precios. Se resisten a que les bajen el nominal, porque empeora su posición frente a otros grupos, pero no a que suban los precios, que reduce a la vez el salario real de todos. El salario nominal es rígido a la baja: $W=W_0$.
\end{itemize}

\subparagraph{ii) Desarrollo}

\begin{itemize}
\item \emph{El empleo lo fija la demanda}: la demanda agregada fija la producción y, con la función de producción, el empleo; el salario real es el que corresponde a ese empleo, $\mathrm{PMg}_L=W_0/P$.
\item \emph{Una caída de la demanda} (gráfico~\ref{fig:keynesiano-shock}):
  \begin{itemize}
  \item en el mercado de bienes (panel b), la demanda agregada pasa de $\mathrm{DA}_0$ a $\mathrm{DA}_1$, y la producción y los precios bajan de $Y_0$ y $P_0$ a $Y_1$ y $P_1$, por debajo de la producción de pleno empleo, $Y_F$;
  \item en el mercado de trabajo (panel a, con el salario nominal en el eje vertical), la demanda de trabajo, el valor de la productividad marginal, pasa de $P_0\cdot\mathrm{PMg}_L$ a $P_1\cdot\mathrm{PMg}_L$;
  \item con el salario nominal fijo en $W_0$, el empleo cae de $L_0$ a $L_1$: el ajuste se hace por cantidades.
  \end{itemize}
\item \emph{El \concepto{paro involuntario}}: es el de quienes trabajarían al salario vigente, o a uno algo menor, y no encuentran empleo. Para \autor{Keynes} lo hay cuando, ante una pequeña subida de los precios respecto del salario nominal, tanto la oferta como la demanda de trabajo a ese salario superarían el empleo existente \parencite[cap.~2]{Keynes1936} \vertema{3.A.28}. Dicho de otro modo: si bajar un poco el salario real por la vía de los precios aumenta el empleo, el paro era involuntario.
\item \emph{La oferta agregada}: mientras haya recursos ociosos tiene pendiente positiva (más demanda sube algo los precios, baja el salario real y aumenta el empleo); cuando se agotan, en $Y_F$, cada subida de precios se traslada a los salarios y la oferta agregada se vuelve vertical, como la neoclásica.
\end{itemize}

\begin{figure}[H]
\centering
\caption{El mercado de trabajo keynesiano ante una caída de la demanda agregada}\label{fig:keynesiano-shock}
\grafico[\textwidth]{keynesiano-shock}
\fuente{elaboración propia a partir de \textcite[cap.~2]{Keynes1936}. OA, oferta agregada.}
\end{figure}

\subparagraph{iii) Implicaciones}

\begin{itemize}
\item \emph{¿Y si bajan los salarios?} \autor{Keynes} no hace descansar el paro en la rigidez del salario: en el capítulo~19 la levanta y muestra que bajar los salarios no lo cura \parencite[cap.~19]{Keynes1936}.
  \begin{itemize}
  \item Una rebaja del salario nominal solo crea empleo si sube la demanda, y para eso tiene que aumentar la inversión: por un tipo de interés más bajo (con precios más bajos sube $M/P$; es el \concepto{efecto \autor{Keynes}}) o por una eficiencia marginal del capital más alta.
  \item Nada garantiza ninguna de las 2 cosas: con trampa de la liquidez el tipo no baja, y la deflación empeora las expectativas y aumenta la carga real de las deudas.
  \item Además, la rebaja quita renta a los trabajadores, que son los que más consumen: se quita a unos lo que se querría dar a otros. Pegatina: «desnudar a un santo para vestir a otro».
  \end{itemize}
  Por eso, para \autor{Keynes}, unos salarios nominales estables dan más estabilidad a la economía que unos salarios flexibles: el salario real no lo fija la negociación salarial, sino el resto del sistema, sobre todo la relación entre la eficiencia marginal del capital y el tipo de interés \parencite[apéndice al cap.~19]{Keynes1936}.
\item \emph{La política}: contra el paro involuntario, política de demanda; bajar los salarios no es la solución.
\item \emph{Recapitulación de la formalización}: en los 3 mercados, \autor{Keynes} cambia la pieza que cerraba el modelo neoclásico. En bienes, la renta, y no el tipo de interés, iguala el ahorro y la inversión; en dinero, el tipo de interés es monetario; en trabajo, el empleo lo fija la demanda y el salario no lo arregla.
\end{itemize}

\begin{itemize}
\item \emph{Los 3 mercados, frente a frente}:
  \begin{itemize}
  \item Cada una de las 3 hojas que siguen resume un mercado en una página: el gráfico, los supuestos, el mecanismo y la conclusión de los neoclásicos y de \autor{Keynes}, lado a lado.
  \item Van en el orden del modelo neoclásico: el trabajo da la oferta agregada (gráfico~\ref{fig:hoja-trabajo}); los bienes (gráfico~\ref{fig:hoja-bienes}) y el dinero (gráfico~\ref{fig:hoja-dinero}), la demanda agregada.
  \item En azul y negrita, los conceptos que introduce \autor{Keynes}; en el recuadro de abajo, la conclusión y lo que implica para la política económica.
  \end{itemize}
\end{itemize}

\begin{figure}[H]
\centering
\caption{Esquema del mercado de trabajo y la oferta agregada: los neoclásicos (\autor{Arthur Cecil Pigou}) frente a \autor{Keynes}}\label{fig:hoja-trabajo}
\grafico[\textwidth]{hoja-trabajo}
\fuente{elaboración propia a partir de \textcite[caps.~2 y 19]{Keynes1936} y \textcite[cap.~1]{Heijdra2017Foundations}. En cada columna, el mercado de trabajo (arriba a la derecha), la función de producción (abajo a la derecha, con la producción hacia abajo) y la oferta agregada (abajo a la izquierda, con el nivel de precios hacia la izquierda).}
\end{figure}

\begin{figure}[H]
\centering
\caption{Esquema del mercado de bienes: los fondos prestables frente a la demanda efectiva}\label{fig:hoja-bienes}
\grafico[\textwidth]{hoja-bienes}
\fuente{elaboración propia a partir de \textcite[caps.~3, 8-12 y 23]{Keynes1936} y \textcite{Fonseca2026Neoclassical}.}
\end{figure}

\begin{figure}[H]
\centering
\caption{Esquema del mercado de dinero y la demanda agregada: la teoría cuantitativa frente a la preferencia por la liquidez}\label{fig:hoja-dinero}
\grafico[\textwidth]{hoja-dinero}
\fuente{elaboración propia a partir de \textcite[caps.~13, 15 y 19]{Keynes1936}, \textcite{Fisher1911Purchasing} y \textcite{Pigou1917Value}.}
\end{figure}

% --------------------------------------------------------------------------
\subsection{Implicaciones de política económica}\label{sec:politica}

\begin{itemize}
\item Si la producción la fija la demanda y nada la lleva sola al pleno empleo, el Estado tiene que sostener la demanda (tabla~\ref{tab:politicas}):
  \begin{itemize}
  \item \emph{la política fiscal} es la más segura: el gasto público entra directamente en la demanda y se multiplica, y con paro no expulsa inversión, porque hay recursos ociosos y el ahorro se crea con la renta. Un déficit en la depresión no es un vicio, sino el remedio, y la obra pública vale aunque su utilidad directa sea escasa, porque lo que importa es el gasto \parencite[cap.~10, apdo.~VI]{Keynes1936};
  \item \emph{la política monetaria} es útil pero insegura, por los 2 eslabones débiles: el tipo baja poco si la preferencia por la liquidez es muy alta, y la inversión responde poco al tipo si las expectativas son malas \vertema{3.A.34}.
  \end{itemize}
\end{itemize}

\begin{table}[H]
\centering
\caption{La política económica en el modelo neoclásico y en el de \autor{Keynes}}\label{tab:politicas}
\small
\begin{tabular}{>{\raggedright\arraybackslash}m{2.6cm}>{\raggedright\arraybackslash}m{5.6cm}>{\raggedright\arraybackslash}m{5.6cm}}
\toprule
\rowcolor{tceeazulpalido} & \textbf{Neoclásicos} & \textbf{\autor{Keynes}}\\
\midrule
Más gasto público & \textcolor{tceerojo}{\textbf{Ineficaz}}: sube el tipo de interés real y expulsa el mismo gasto privado; la producción no cambia & \textcolor{tceeverde}{\textbf{Eficaz}}: la demanda sube y se multiplica; con recursos ociosos suben la producción y el empleo\\
\addlinespace[0.3em]
Más dinero & \textcolor{tceerojo}{\textbf{Ineficaz}}: solo sube los precios en la misma proporción (neutralidad) & \textcolor{tceedorado}{\textbf{Insegura}}: baja el tipo, sube la inversión y la renta, salvo trampa de la liquidez o inversión insensible\\
\addlinespace[0.3em]
Bajar los salarios & \textcolor{tceeverde}{\textbf{Eficaz}}: vacía el mercado de trabajo y devuelve el pleno empleo & \textcolor{tceerojo}{\textbf{Ineficaz}}: no asegura más empleo y puede agravar la depresión (cap.~19)\\
\addlinespace[0.3em]
Recomendación & \emph{Laissez-faire} y presupuesto equilibrado & Gestión de la demanda, sobre todo fiscal; socialización de la inversión (cap.~24)\\
\bottomrule
\end{tabular}
\fuente{elaboración propia a partir de \textcite{Keynes1936} y \textcite[cap.~1]{Heijdra2017Foundations}. En verde, lo eficaz; en rojo, lo ineficaz; en dorado, lo que depende de las circunstancias.}
\end{table}

\begin{itemize}
\item \emph{Más allá de la coyuntura} \parencite[cap.~24]{Keynes1936}:
  \begin{itemize}
  \item un tipo de interés bajo de forma permanente, que lleve a la «eutanasia del rentista»: el capital dejaría de ser escaso y de rendir solo por serlo;
  \item «una socialización bastante completa de la inversión» como único medio de acercarse al pleno empleo \parencite[cap.~24; traducción propia]{Keynes1936}, que no quiere decir nacionalizar la propiedad, sino que el Estado asegure el volumen de inversión;
  \item una redistribución de la renta con impuestos progresivos, que sube la propensión a consumir de la sociedad.
  \end{itemize}
\item \emph{La oferta agregada}: con la oferta agregada creciente, la política de demanda mueve la producción y el empleo mientras haya recursos ociosos; cuando se llega a $Y_F$, solo mueve los precios \vertema{3.A.32}.
\item \emph{La inflación}: una expansión puntual sube los precios una sola vez; para mantener el empleo alto habría que repetirla, y eso ya es inflación. A la teoría de \autor{Keynes} le falta una teoría de la inflación: la curva de \autor{Phillips} y su crítica la darán después \vertema{3.A.5} \vertema{3.A.28}.
\end{itemize}

% --------------------------------------------------------------------------
\subsection{Valoración: las lecturas de la \emph{Teoría general}}\label{sec:lecturas}

\begin{itemize}
\item \emph{Virtudes}:
  \begin{itemize}
  \item \emph{explica el paro masivo} de los años 30 como un fallo del sistema, y no del mercado de trabajo, y da una receta para combatirlo;
  \item \emph{funda la macroeconomía}: un análisis de la economía entera, simplificado y orientado a la política \parencite{DeVroey2016History};
  \item \emph{rompe la dicotomía clásica} a corto plazo: el dinero y las expectativas tienen efectos reales.
  \end{itemize}
  Su éxito entre los economistas jóvenes fue inmediato: para \textcite[pág.~187; traducción propia]{Samuelson1946Lord}, la \emph{Teoría general} «contagió a la mayoría de los economistas de menos de 35 años con la virulencia inesperada de una enfermedad que ataca por primera vez a una tribu aislada de los mares del Sur y la diezma».
\item \emph{Críticas}:
  \begin{itemize}
  \item \emph{la demostración es incompleta}: \autor{Keynes} quería probar que puede haber paro involuntario sin rigideces, y no lo consiguió con el instrumental de su tiempo; el libro se leyó como una teoría de salarios rígidos, justo lo que quería evitar \parencite{DeVroey2016History};
  \item \emph{es estático y de corto plazo}: no tiene una teoría de la inflación ni del crecimiento, que intentarán \autor{Harrod} y los postkeynesianos (\ref{sec:postkeynesianos});
  \item \emph{no tiene fundamentos microeconómicos} explícitos para la demanda efectiva: será el programa de los teóricos del desequilibrio (\ref{sec:desequilibrio}) y, después, de la nueva economía keynesiana \vertema{3.A.7}.
  \end{itemize}
\item \emph{La lectura de \autor{Hicks}: el IS-LM} \vertema{3.A.5}:
  \begin{itemize}
  \item \textcite{Hicks1937Keynes} reduce la \emph{Teoría general} a 2 curvas: la IS, los pares de renta y tipo de interés que equilibran el mercado de bienes (el multiplicador \eqref{eq:multiplicador} con el tipo variable), y la LM, los que equilibran el de dinero \eqref{eq:liquidez}. La renta y el tipo se determinan a la vez.
  \item \textcite{Modigliani1944Liquidity} muestra que, con salarios flexibles y sin trampa de la liquidez, el modelo vuelve al pleno empleo: el paro keynesiano pasa a depender de la rigidez del salario nominal.
  \item Es el modelo de la \emph{síntesis neoclásica}, que lo completará con la curva de \autor{Phillips}. Para la síntesis, \autor{Keynes} queda como el caso particular de los salarios rígidos, no como una «teoría general».
  \end{itemize}
\item \emph{¿Traducción o traición?} Para \textcite{Patinkin1987Keynes}, la \emph{Teoría general} hablaba con la voz de \autor{Marshall}, pero razonaba con el sistema de ecuaciones interdependientes de \autor{Walras}; la síntesis se quedó con lo segundo. Para unos, la síntesis tradujo a \autor{Keynes}; para otros, lo traicionó.
\item \emph{Recapitulación}: la \emph{Teoría general} da la vuelta a la causalidad neoclásica (la demanda fija la producción y el empleo; el dinero, el tipo de interés; la renta iguala el ahorro y la inversión) y deja abierta la pregunta de cuál es exactamente su modelo. Una vez vista la lectura de la síntesis, pasamos a las que se enfrentaron a ella.
\end{itemize}

% ==========================================================================
\section{La economía postkeynesiana y el desequilibrio}\label{sec:herederos}
\minutos{10}

\begin{itemize}
\item Una vez vista la lectura de la síntesis, pasamos a las 2 que se enfrentaron a ella.
\item \emph{Conclusión anticipada}: los postkeynesianos rechazan el instrumental neoclásico y llevan la demanda efectiva y la incertidumbre al largo plazo, a la distribución y a las finanzas. Los keynesianos del desequilibrio hacen lo contrario: aceptan ese instrumental y buscan en él los fundamentos de \autor{Keynes}, en una economía en la que los precios no vacían los mercados y los agentes responden a las cantidades. Los 2 coinciden en que el paro no se explica solo con un salario que no baja.
\end{itemize}

% --------------------------------------------------------------------------
\subsection{La economía postkeynesiana: la escuela de Cambridge}\label{sec:postkeynesianos}

\paragraph{a) Caracterización según \autor{Schumpeter}}

\begin{itemize}
\item \emph{Líder}: \autor{Joan Robinson}, con \autor{Nicholas Kaldor}. El núcleo es el «círculo» que discutió con \autor{Keynes} el \emph{Treatise} y la \emph{Teoría general}: \autor{Richard Kahn}, que había formulado el multiplicador \parencite{Kahn1931Relation}, \autor{Joan Robinson} y \autor{Piero Sraffa}. Se les unen después \autor{Michał Kalecki}, su precursor, y \autor{Luigi Pasinetti}; en Estados Unidos, \autor{Hyman Minsky}, \autor{Paul Davidson} y \autor{Sidney Weintraub} \parencite{King2002History}.
\item \emph{Su reacción a la síntesis}: para \autor{Robinson}, el IS-LM era un «keynesianismo bastardo» \parencite{Robinson1962Review}: había devuelto a \autor{Keynes} al equilibrio neoclásico y olvidado sus 2 lecciones, el tiempo histórico (el futuro es incierto y las decisiones no se pueden deshacer) y la demanda efectiva también a largo plazo.
\item \emph{Objeto}, con 4 ideas comunes \parencite{Lavoie2014PostKeynesian}:
  \begin{itemize}
  \item la \emph{demanda efectiva} manda también a largo plazo: la inversión determina el ahorro, y no al revés;
  \item la \emph{incertidumbre fundamental}: ninguna distribución de probabilidad la reduce a riesgo;
  \item la \emph{distribución} de la renta depende de las instituciones y del poder de mercado, no de las productividades marginales, y a su vez afecta a la demanda;
  \item el \emph{dinero} es endógeno y el sistema financiero, una fuente de inestabilidad.
  \end{itemize}
\item \emph{Método}: tiempo histórico frente a equilibrio, clases sociales (trabajadores y capitalistas) en lugar de un agente representativo, precios fijados con un margen sobre los costes y rechazo de la función de producción agregada.
\end{itemize}

\paragraph{b) \autor{Kalecki}: demanda efectiva y distribución}

\begin{itemize}
\item \emph{Idea}: \autor{Michał Kalecki} llegó a la demanda efectiva por su cuenta, desde los esquemas de reproducción de \autor{Marx}, en un ensayo sobre el ciclo publicado en polaco en 1933, antes que la \emph{Teoría general} \parencite{Kalecki1933Proba,Kalecki1971Selected}.
\item \emph{Los beneficios}: en una economía cerrada, sin sector público y en la que los trabajadores gastan todo lo que ganan, los beneficios $B$ son iguales a la inversión más el consumo de los capitalistas, $C_{\mathrm{B}}$:
  \begin{equation}
  B=I+C_{\mathrm{B}}. \label{eq:beneficios-kalecki}
  \end{equation}
  La causalidad va del gasto a los beneficios: los capitalistas deciden cuánto gastan, no cuánto ganan. \textcite{Kaldor1956Alternative} lo resumió así: los capitalistas ganan lo que gastan y los trabajadores gastan lo que ganan.
\item \emph{El \concepto{grado de monopolio}}: las empresas fijan el precio con un margen $\mu$ sobre el coste directo (salarios y materias primas). Si $j$ es la relación entre el coste de las materias primas y los salarios, la parte de los salarios en el valor añadido es
  \begin{equation}
  \frac{W\,L}{P\,Y}=\frac{1}{1+\mu\,(1+j)}, \label{eq:kalecki}
  \end{equation}
  y depende del poder de mercado de las empresas (la concentración, la publicidad, la fuerza de los sindicatos), no de la productividad marginal \parencite{Kalecki1938Determinants,Kalecki1954Theory}. El margen está ligado al índice de \autor{Lerner}, $\mathcal{L}=\mu/(1+\mu)$ \vertema{3.A.17}.\footnote{Con ingresos $(1+\mu)$ veces el coste directo, el valor añadido es $(1+\mu)\,W\,L+\mu$ veces el coste de las materias primas; al dividir la masa salarial por esa cantidad sale \eqref{eq:kalecki}. \autor{Kalecki} escribe $k=1+\mu$, pero aquí $k$ es la de la demanda de dinero.}
\item \emph{Implicaciones}: si los trabajadores consumen más que los capitalistas, subir los salarios reales puede aumentar la demanda, el empleo y hasta los beneficios totales: es la «paradoja de los costes». Que una economía crezca más con salarios altos o con beneficios altos depende de qué pese más, el consumo o la inversión que se resiente con márgenes más bajos \parencite{BhaduriMarglin1990Unemployment} \vertema{3.A.45}.
\item \emph{El ciclo político}: \textcite{Kalecki1943Political} explica por qué los gobiernos no mantienen el pleno empleo aunque sepan cómo: el pleno empleo duradero da poder a los trabajadores y erosiona la disciplina en las fábricas, y los empresarios presionan para volver a la austeridad. Es un antecedente de los ciclos político-económicos \vertema{3.A.23}.
\end{itemize}

\paragraph{c) Crecimiento y distribución: \autor{Harrod}, \autor{Kaldor} y \autor{Pasinetti}}

\begin{itemize}
\item \autor{Harrod} lleva la demanda efectiva al crecimiento \parencite{Harrod1939}: la tasa de crecimiento \emph{garantizada}, $s/v$, la que cumple las expectativas de los empresarios, solo coincide por casualidad con la \emph{natural}, la del pleno empleo, y si la economía se aparta de la garantizada, se aleja cada vez más: es el «filo de la navaja» \vertema{3.A.43}.
\item \emph{La salida de Cambridge (Inglaterra)}: \autor{Solow} igualará las 2 tasas haciendo variable la relación capital-producto, $v$; \textcite{Kaldor1956Alternative} lo hace con la \emph{distribución}. Si los capitalistas ahorran una parte $s_{\mathrm{B}}$ de los beneficios y los trabajadores una parte $s_{\mathrm{L}}<s_{\mathrm{B}}$ de los salarios, la igualdad entre el ahorro y la inversión, $s_{\mathrm{B}}\,B+s_{\mathrm{L}}\,(Y-B)=I$, da
  \begin{equation}
  \frac{B}{Y}=\frac{1}{s_{\mathrm{B}}-s_{\mathrm{L}}}\,\frac{I}{Y}-\frac{s_{\mathrm{L}}}{s_{\mathrm{B}}-s_{\mathrm{L}}}. \label{eq:kaldor}
  \end{equation}
  \begin{itemize}
  \item La inversión fija la parte de los beneficios, y no la productividad marginal del capital: si sube la inversión, sube la parte de los beneficios y, con ella, el ahorro que la financia. La tasa de ahorro deja de ser un dato.
  \end{itemize}
\item \textcite{Pasinetti1962Rate} añade que los trabajadores que ahorran acaban teniendo capital y cobrando beneficios. Aun así, en una senda de crecimiento equilibrado la tasa de beneficio solo depende del crecimiento y de lo que ahorran los capitalistas:
  \begin{equation}
  \frac{B}{K}=\frac{\gamma_Y}{s_{\mathrm{B}}}. \label{eq:pasinetti}
  \end{equation}
  Es la \emph{ecuación de Cambridge}: no depende de la tecnología ni de lo que ahorren los trabajadores \vertema{3.A.43} \vertema{3.A.45}.\footnote{Notación local: $s_{\mathrm{B}}$ y $s_{\mathrm{L}}$ son las propensiones a ahorrar de los beneficios y de los salarios (\autor{Kaldor} escribe $s_p$ y $s_w$), y la tasa de beneficio se escribe $B/K$ porque $r$ es en el temario el tipo de interés real.}
\end{itemize}

\paragraph{d) La controversia del capital}

\begin{itemize}
\item \emph{La pregunta}: si el capital son máquinas distintas, ¿en qué unidades se suma para meterlo en una función de producción agregada? \textcite{Robinson1953Production} vio la trampa: solo puede sumarse en valor, y ese valor depende del tipo de beneficio, que es lo que la teoría quería explicar con la productividad marginal del capital. Pegatina: «no se pueden sumar peras con manzanas».
\item \emph{El resultado}: \textcite{Sraffa1960Production} construye el contraejemplo. Cada técnica tiene una \emph{frontera salario-beneficio}: cuanto más se paga al trabajo, menos queda para el capital (gráfico~\ref{fig:retorno-tecnicas}). Si 2 fronteras se cortan 2 veces, la técnica $\alpha$ es la más rentable con tipos bajos, la $\beta$ con tipos intermedios y la $\alpha$ otra vez con tipos altos: es el \concepto{retorno de técnicas} (\emph{reswitching}). Entonces no puede ordenarse una técnica como «más intensiva en capital» que otra, y el capital por trabajador puede subir cuando sube el tipo de beneficio (\emph{reversión del capital}).
\item \emph{El final}: el intento de \textcite{Levhari1965Nonswitching} de probar que eso no ocurre en la economía en su conjunto resultó erróneo \parencite{LevhariSamuelson1966Nonswitching}, y \textcite{Samuelson1966Summing} lo admitió desde Cambridge (Massachusetts) \parencite{CohenHarcourt2003Whatever}.
\item \emph{Lo que significa para \autor{Keynes} y la distribución}: si no hay una cantidad de capital que medir antes de conocer el tipo de beneficio, la distribución no puede explicarse por las productividades marginales; la explican la acumulación (\autor{Kaldor}) y el poder de negociación (\autor{Kalecki}). La macroeconomía, con todo, siguió usando la función agregada como aproximación útil \vertema{3.A.11} \vertema{3.A.43}.
\end{itemize}

\begin{figure}[H]
\centering
\caption{El retorno de técnicas: 2 fronteras salario-beneficio que se cortan 2 veces}\label{fig:retorno-tecnicas}
\grafico[0.55\textwidth]{retorno-tecnicas}
\fuente{elaboración propia a partir de \textcite{Sraffa1960Production} y \textcite{Samuelson1966Summing}. En el eje horizontal, el tipo de beneficio (que en este debate coincide con el de interés); en el vertical, el salario real.}
\end{figure}

\paragraph{e) \autor{Minsky}: la hipótesis de la inestabilidad financiera}

\begin{itemize}
\item \emph{Idea}: \autor{Hyman Minsky} lee a \autor{Keynes} como el teórico de una economía con bancos y deudas, en la que la inestabilidad nace dentro del sistema financiero \parencite{Minsky1975Keynes,Minsky1986Stabilizing}.
\item \emph{Las 3 posiciones financieras} \parencite{Minsky1992Financial}:
  \begin{itemize}
  \item de \emph{cobertura}: los ingresos esperados pagan los intereses y el principal;
  \item \emph{especulativa}: pagan los intereses, pero el principal hay que refinanciarlo;
  \item \emph{Ponzi}: no pagan ni los intereses; hay que endeudarse más o vender activos, y solo funciona si sus precios suben.
  \end{itemize}
\item \emph{La dinámica}: en una expansión larga, los éxitos animan a endeudarse y los prestamistas relajan sus criterios; las posiciones pasan de cobertura a especulativas y de especulativas a Ponzi. La estabilidad es desestabilizadora: basta una subida de tipos o una caída de los ingresos para que las ventas forzosas hundan los precios de los activos y la crisis se propague. Pegatina: «después de la calma viene la tempestad».
\item \emph{Implicaciones}: un «gran gobierno», cuyo déficit sostiene los beneficios en la recesión, y un «gran banco», el banco central como prestamista de última instancia, frenan la caída; pero, al hacerlo, animan a asumir más riesgo, y la regulación financiera tiene que ir por delante \parencite{Minsky1986Stabilizing} \vertema{3.A.42}.
\item \emph{El dinero endógeno}: para los postkeynesianos, el dinero no lo fija el banco central, como en la \emph{Teoría general}, sino que lo crean los bancos al conceder créditos; el banco central fija el tipo de interés y da las reservas que se le piden, y la oferta de dinero es «horizontal» a ese tipo \parencite{Kaldor1982Scourge,Moore1988Horizontalists} \vertema{3.A.37}.\footnote{El postkeynesianismo estadounidense, con \textcite{Davidson1972Money} y \autor{Weintraub}, insistió en el dinero y en los contratos nominales como respuesta a la incertidumbre; \textcite{KregelEichner1975Essay} lo presentaron como un nuevo paradigma.}
\end{itemize}

\paragraph{f) Evidencia empírica}

\begin{itemize}
\item La crisis de 2007-2008 se leyó como un «momento \autor{Minsky}» \parencite{Cassidy2008Minsky}: la deuda del sector privado no financiero de Estados Unidos pasó del 125~\% del PIB en 1990 al 171~\% a finales de 2007, y la de los hogares, del 61~\% al 98~\%, con una parte creciente que solo se sostenía si la vivienda seguía subiendo (gráfico~\ref{fig:deuda-privada}). El máximo, 172,0~\%, llegó en el tercer trimestre de 2008; a finales de 2025 era del 140,3~\% \parencite{BIS2026Credito} \vertema{3.A.42}.
\item El Banco de Inglaterra describe hoy la creación de dinero como los postkeynesianos: la mayor parte del dinero son depósitos que crean los bancos al prestar \parencite{McLeayRadiaThomas2014Money}.
\end{itemize}

\begin{figure}[H]
\centering
\caption{Crédito al sector privado no financiero de Estados Unidos, 1980-2025 (en porcentaje del PIB)}\label{fig:deuda-privada}
\grafico[0.6\textwidth]{eeuu-deuda-privada}
\fuente{elaboración propia a partir de \textcite{BIS2026Credito}; datos trimestrales, último dato del cuarto trimestre de 2025. En gris, las recesiones.}
\end{figure}

\paragraph{g) Valoración}

\begin{itemize}
\item \emph{Virtudes}: tomaron en serio la incertidumbre, la distribución y las finanzas, que la macroeconomía dominante dejó de lado hasta 2008, y anticiparon el dinero endógeno, hoy aceptado.
\item \emph{Críticas}: no formaron un modelo común ni un programa con contrastes empíricos sistemáticos; su rechazo del instrumental neoclásico los dejó fuera de la corriente principal y de la política; y la controversia del capital, que ganaron en la lógica, no cambió la práctica.
\item \emph{Recapitulación}: los postkeynesianos se quedan con la demanda efectiva y la incertidumbre y rechazan el instrumental neoclásico. Otros economistas hicieron lo contrario: aceptarlo y buscar en él lo que le faltaba a \autor{Keynes}.
\end{itemize}

% --------------------------------------------------------------------------
\subsection{Los keynesianos del desequilibrio}\label{sec:desequilibrio}

\paragraph{Caracterización según \autor{Schumpeter}}

\begin{itemize}
\item \emph{Líderes}: \autor{Robert Clower} y \autor{Axel Leijonhufvud}, que a finales de los 60 sostienen que la \emph{Teoría general} es una teoría de lo que pasa cuando los precios \emph{no} vacían los mercados, y que el IS-LM, un modelo de equilibrio, la había perdido por el camino.
\item \emph{Seguidores}: \autor{Don Patinkin} como antecedente; los modelos de \autor{Robert Barro} y \autor{Herschel Grossman}, \autor{Jean-Pascal Bénassy}, \autor{Jacques Drèze} y \autor{Edmond Malinvaud}. Es una escuela de los 70, sobre todo europea, que después se dispersa \parencite{DeVroey2016History}.
\item \emph{Objeto}: los equilibrios sin vaciado de mercado, llamados \concepto{equilibrios no walrasianos} o \emph{con racionamiento}:
  \begin{itemize}
  \item los precios se ajustan despacio, tanto que a corto plazo pueden tomarse como dados;
  \item se comercia a precios que no son los de equilibrio: en cada mercado se intercambia lo que quiere el \emph{lado corto}, y el otro lado queda \emph{racionado};
  \item el racionamiento en un mercado se contagia a los demás, y la economía puede quedarse con paro involuntario. Es un problema de \emph{coordinación}: las empresas no contratan porque no venden, y los hogares no compran porque no tienen empleo.
  \end{itemize}
\item \emph{Método}: hogares y empresas que maximizan, como en el equilibrio general, pero sin subastador walrasiano y con restricciones de cantidad además de las de precios.
\end{itemize}

\paragraph{\autor{Patinkin} (1956): el paro como ajuste lento}

\subparagraph{a) Idea}

\begin{itemize}
\item En \emph{Money, Interest, and Prices}, \textcite{Patinkin1956Money} (2.ª ed., 1965) integra el dinero en el equilibrio general de \autor{Walras}: con los saldos reales en la demanda de bienes, el dinero fija el nivel de precios sin romper la teoría del valor \vertema{3.A.21}. Y quiere hacer sitio en ese equilibrio general al paro involuntario de \autor{Keynes}.
\end{itemize}

\subparagraph{b) Modelo}

\begin{itemize}
\item \emph{i) Supuestos}: 4 mercados (trabajo, bienes, bonos y dinero), con la atención en los 2 primeros; precios y salarios flexibles, pero \emph{lentos}.
\item \emph{ii) Desarrollo} (gráfico~\ref{fig:patinkin}, que redibuja los de \textcite[cap.~3]{DeVroey2016History}):
  \begin{itemize}
  \item en el trabajo, la demanda sale de $\mathrm{PMg}_L(L)=w$ y la oferta es $L^S(w)$; en el equilibrio, $(L_0,w_0)$ y $Y_0=F(L_0)$: el punto $A$ del panel a y el $B$ del panel b, donde el gasto planeado, $E_0$, iguala la producción. La demanda de bienes incluye el \concepto{efecto saldos reales} (o efecto \autor{Pigou}-\autor{Patinkin}): si baja $P$, sube la riqueza real y el consumo,
    \begin{equation}
    Y^D=C\Big(Y,\,i,\,\frac{M}{P}\Big)+I(i)+\overline{G},\qquad \frac{\partial C}{\partial(M/P)}>0; \label{eq:patinkin}
    \end{equation}
  \item \emph{la perturbación}: cae la demanda de bienes. Como los precios bajan despacio, las empresas no venden lo que producían, acumulan existencias y acaban produciendo menos. Queda un \emph{exceso de oferta de bienes}: producirían más si se lo compraran (de $C$ a $B$), pero solo venden $Y_1$ (el punto $G$);
  \item necesitan menos trabajo y, al mismo salario real (si precios y salarios bajan en la misma proporción), contratan menos: el empleo queda en el punto $K$, \emph{fuera} de su curva de demanda de trabajo y de la curva de oferta de los trabajadores. Ni unos ni otros eligen: la falta de demanda de bienes los fuerza.
  \end{itemize}
\item \emph{iii) Implicaciones}:
  \begin{itemize}
  \item hay \emph{paro involuntario sin que el salario real esté por encima del de equilibrio}: el paro nace en el mercado de bienes y se traslada al de trabajo, que es la tesis de \autor{Keynes};
  \item pero no es un equilibrio: con exceso de oferta de bienes y de trabajo, precios y salarios siguen bajando; el efecto saldos reales sube la demanda ($E_1\to E_2\to E_0$) y la economía vuelve, despacio, de $K$ a $A$ y de $G$ a $B$: al pleno empleo;
  \item la política de demanda sirve para \emph{acortar} un ajuste que llegaría solo.
  \end{itemize}
\end{itemize}

\begin{figure}[H]
\centering
\caption{\autor{Patinkin}: el paro involuntario como ajuste lento tras una caída de la demanda de bienes}\label{fig:patinkin}
\grafico[0.95\textwidth]{patinkin-desequilibrio}
\fuente{elaboración propia a partir de \textcite[figs.~3.1 y 3.2]{DeVroey2016History} y \textcite{Patinkin1956Money}. $w$, salario real; $E$, gasto planeado; $Y=S(w_0,K_0)$, la producción que las empresas quieren ofrecer con el salario y el capital dados.}
\end{figure}

\subparagraph{c) Valoración}

\begin{itemize}
\item \emph{Virtudes}: primera explicación del paro involuntario dentro del equilibrio general sin un salario demasiado alto; para \autor{Patinkin}, la verdadera síntesis de \autor{Keynes} y \autor{Walras}, frente a la del IS-LM, que reducía a \autor{Keynes} al caso del salario rígido o de la trampa de la liquidez.
\item \emph{Críticas}: el paro es solo transitorio, no el «equilibrio con paro» que buscaba \autor{Keynes}; da por hecho que el ajuste es estable; y el efecto saldos reales es pequeño y lento \parencite{DeVroey2016History}. Para \autor{Patinkin}, si el paro se debía solo a salarios rígidos, \autor{Keynes} no aportaba nada nuevo; \textcite{Lucas1980Methods} vio en su libro la versión más acabada de la síntesis.
\end{itemize}

\paragraph{\autor{Clower} (1965): la decisión dual}

\subparagraph{a) Idea}

\begin{itemize}
\item En «The Keynesian Counter-Revolution», \textcite{Clower1965Keynesian} busca lo nuevo de \autor{Keynes} en la función de consumo: el consumo depende de la renta \emph{realizada}, algo que la teoría walrasiana no puede recoger, porque en ella la demanda depende solo de los precios. De ahí su dilema: o la ley de \autor{Walras} es incompatible con la economía keynesiana, «o \autor{Keynes} no tenía nada fundamentalmente nuevo que añadir a la teoría ortodoxa» \parencite[traducción propia]{Clower1965Keynesian}.
\end{itemize}

\subparagraph{b) Modelo}

\begin{itemize}
\item \emph{i) Supuestos}: no hay subastador; se comercia a los precios vigentes aunque no vacíen los mercados; los hogares venden trabajo y compran bienes.
\item \emph{ii) Desarrollo} (gráfico~\ref{fig:clower}):
  \begin{itemize}
  \item la \concepto{demanda nocional} es la walrasiana, la que el hogar planea si puede vender todo el trabajo que quiere a los precios vigentes: $C^{\mathrm{noc}}$, con $L^S(w)$;
  \item la \concepto{demanda efectiva}: si no puede vender más que $\overline{L}<L^S(w)$, su demanda de bienes depende de la renta que de verdad obtiene, $C^{\mathrm{ef}}=C(w\,\overline{L})<C^{\mathrm{noc}}$. Es la \concepto{hipótesis de la decisión dual}: el hogar decide primero cuánto trabajo vende y, con la renta que consigue, cuánto compra. Pegatina: «de donde no hay, no se puede sacar».\footnote{«Demanda efectiva» tiene así 3 sentidos en el temario: el punto de demanda efectiva de \autor{Keynes} (capítulo~3), la demanda condicionada por las restricciones de cantidad de \autor{Clower} y la curva de cuota de mercado de \autor{Chamberlin} \vertema{3.A.18}.}
  \item \emph{La ley de \autor{Walras}} \vertema{3.A.21}: con las demandas nocionales, el valor de los excesos de demanda de los $n$ mercados es siempre cero,
    \begin{equation}
    \sum_{k=1}^{n}p_k\,z_k^{\mathrm{noc}}(\vect{p})\equiv 0. \label{eq:walras-nocional}
    \end{equation}
    Con el trabajo racionado, el hogar gasta en bienes y dinero lo que de verdad ingresa, y el valor de sus excesos de demanda efectivos de bienes y de dinero es cero. Pero el exceso de oferta de trabajo sigue ahí; como el empleo lo fija el lado corto, las empresas contratan $L^D=\overline{L}$:
    \begin{equation}
    \sum_{k\neq L}p_k\,z_k^{\mathrm{ef}}+W\,\big(L^D-L^S\big)=W\,\big(\overline{L}-L^S\big)<0. \label{eq:walras-efectiva}
    \end{equation}
    Al paro no le corresponde, al otro lado, un exceso de demanda de bienes: los parados querrían comprar, pero su deseo no es una demanda efectiva.
  \item \emph{Un ejemplo}: un consultor sin clientes querría beber champán, y el bodeguero querría vender más champán y contratar consultores para crecer. Los 2 saldrían ganando con el intercambio, pero el deseo del consultor no llega al mercado: sin renta, no es demanda, y el precio no transmite la señal que los coordinaría.
  \end{itemize}
\item \emph{iii) Implicaciones}: el paro no envía ninguna señal al mercado de bienes; nada empuja a las empresas a producir más. La economía puede quedarse en un equilibrio con paro, un fallo de coordinación que la política de demanda resuelve; el multiplicador es la forma de ese contagio entre mercados.
\end{itemize}

\begin{figure}[H]
\centering
\caption{La decisión dual de \autor{Clower}: demanda nocional y demanda efectiva de un hogar racionado en el empleo}\label{fig:clower}
\grafico[0.55\textwidth]{clower-decision-dual}
\fuente{elaboración propia a partir de \textcite{Clower1965Keynesian} y \textcite{BarroGrossman1971General}.}
\end{figure}

\subparagraph{c) Extensiones}

\begin{itemize}
\item \textcite{Clower1967Reconsideration} añade que, en una economía monetaria, los bienes se compran con dinero: el dinero tiene que estar antes que la compra \vertema{3.A.21}.
\end{itemize}

\subparagraph{d) Evidencia empírica}

\begin{itemize}
\item El consumo responde a la renta corriente más de lo que predice la renta permanente, como cabe esperar si una parte de los hogares no puede gastar lo que no ingresa \vertema{3.A.33}.
\end{itemize}

\subparagraph{e) Valoración}

\begin{itemize}
\item \emph{Virtudes}: da fundamento microeconómico a la función de consumo keynesiana y al multiplicador.
\item \emph{Críticas}: si todo se mide en términos efectivos (con $\overline{L}$ en lugar de $L^S$), la ley de \autor{Walras} vuelve a cumplirse; y la decisión dual explica qué pasa \emph{una vez} que los precios están mal, pero no por qué no se ajustan. Para \autor{Clower}, aun así, la síntesis había reducido a una cuestión de salarios rígidos una teoría revolucionaria.
\end{itemize}

\paragraph{\autor{Leijonhufvud} (1968): la economía de \autor{Keynes} frente a la economía keynesiana}

\subparagraph{a) Idea}

\begin{itemize}
\item En \emph{On Keynesian Economics and the Economics of Keynes}, \textcite{Leijonhufvud1968Keynesian} separa la «economía keynesiana», la de la síntesis, de la «economía de \autor{Keynes}», la del libro, que para él era una teoría coherente que \autor{Keynes} no supo modelizar.
\end{itemize}

\subparagraph{b) Modelo: sus 4 tesis}

\begin{itemize}
\item \emph{La velocidad de ajuste}: en \autor{Keynes}, a corto plazo las cantidades se ajustan antes que los precios; es el orden inverso al de \autor{Marshall}. No hace falta suponer precios rígidos: basta con que se muevan más despacio que las cantidades.
\item \emph{La información}: sin subastador, nadie conoce los precios de equilibrio; mientras los buscan, vendedores y trabajadores mantienen sus precios y lo que se ajusta es lo que se vende y se emplea.
\item \emph{El multiplicador}: la caída de las ventas reduce la renta y, con ella, la demanda efectiva de otros (la decisión dual): un proceso que amplifica la perturbación.
\item \emph{El tipo de interés}: el problema de fondo no es el salario, sino un tipo de interés a largo plazo que no baja lo bastante cuando cae la inversión, porque los inversores esperan que vuelva a su nivel «normal»: el tipo de mercado queda por encima del natural de \autor{Wicksell}. El ahorro y la inversión no se coordinan en el tiempo, y bajar los salarios no ataca la causa.
\end{itemize}
\begin{itemize}
\item \emph{Implicaciones}: el paro keynesiano es un fallo de coordinación por falta de información, y no la consecuencia de un salario rígido.\footnote{En «Effective Demand Failures», \textcite{Leijonhufvud1973Effective} propone el \emph{pasillo} (\emph{corridor}): dentro de una banda alrededor del equilibrio, los colchones de la economía (existencias, dinero, crédito) absorben las perturbaciones y los precios devuelven la economía al equilibrio; fuera de ella, los colchones se agotan, el gasto pasa a depender de la renta corriente y el multiplicador amplifica la caída. Los neoclásicos tendrían razón para las perturbaciones pequeñas y \autor{Keynes}, para las grandes, como la de 1929.}
\end{itemize}

\subparagraph{c) Valoración}

\begin{itemize}
\item \emph{Virtudes}: devolvió a \autor{Keynes} el problema de la coordinación y la información, y explicó por qué la síntesis lo había perdido.
\item \emph{Críticas}: es una interpretación discutida de \autor{Keynes}, no un modelo formal; la parte formal la harán \autor{Barro} y \autor{Grossman} y \autor{Malinvaud}.
\end{itemize}

\paragraph{\autor{Barro} y \autor{Grossman} (1971): el modelo de desequilibrio general}

\subparagraph{a) Idea}

\begin{itemize}
\item \textcite{BarroGrossman1971General} juntan la demanda efectiva de bienes de \autor{Clower} y la demanda de trabajo racionada de \autor{Patinkin} en un modelo de equilibrio general con precios y salarios fijos.
\end{itemize}

\subparagraph{b) Modelo}

\begin{itemize}
\item \emph{i) Supuestos}: el nivel de precios $P$ y el salario nominal $W$ están dados; empresas competitivas con $Y=F(L)$; en cada mercado manda el lado corto.
\item \emph{ii) Desarrollo}, con exceso de oferta de bienes y de trabajo:
  \begin{itemize}
  \item las empresas, que no venden todo lo que querrían, solo contratan el trabajo que necesitan para producir lo que venden: su \emph{demanda efectiva de trabajo}, $L^{\mathrm{ef}}=F^{-1}(Y)$, es menor que la nocional, la de $\mathrm{PMg}_L=W/P$;
  \item los hogares, racionados en el empleo, gastan según la renta que obtienen: $C^{\mathrm{ef}}=C\big(\tfrac{W}{P}\,L^{\mathrm{ef}},\,\tfrac{M}{P}\big)$;
  \item la producción es la demandada, $Y=C^{\mathrm{ef}}+\overline{I}+\overline{G}$: cada mercado limita al otro y el resultado es el multiplicador keynesiano.
  \end{itemize}
  En el caso opuesto, con exceso de demanda de bienes y de trabajo, los hogares no encuentran qué comprar, deciden trabajar menos y aparece un multiplicador de oferta.
\item \emph{iii) Implicaciones}: puede haber paro involuntario con el salario real «correcto»; bajarlo no ayuda si el problema es la demanda.
\item \emph{Qué añade frente a \autor{Patinkin}} (gráfico~\ref{fig:racionamiento}, con los mismos ejes que el gráfico~\ref{fig:patinkin}):
  \begin{itemize}
  \item las empresas que no venden solo demandan el trabajo necesario para producir $Y_1$: la demanda efectiva de trabajo se vuelve vertical en $N_1$, y el empleo queda en $K$ con el salario real del equilibrio walrasiano;
  \item los hogares racionados gastan según la renta que obtienen: la demanda efectiva de bienes, $E^{\mathrm{ef}}$, queda por debajo de la nocional, y la economía se queda en $G$;
  \item con $P$ y $W$ fijos no hay efecto saldos reales: $K$ y $G$ ya no son un tránsito hacia $A$ y $B$, sino un \emph{equilibrio con racionamiento}, y solo un aumento del gasto autónomo, con su multiplicador, saca de él a la economía.
  \end{itemize}
\end{itemize}

\begin{figure}[H]
\centering
\caption{\autor{Clower} y \autor{Barro} y \autor{Grossman}: el paro keynesiano como equilibrio con racionamiento}\label{fig:racionamiento}
\grafico[0.95\textwidth]{desequilibrio-racionamiento}
\fuente{elaboración propia a partir de \textcite{Clower1965Keynesian}, \textcite{BarroGrossman1971General} y \textcite[cap.~3]{DeVroey2016History}. $E^{\mathrm{noc}}$ y $E^{\mathrm{ef}}$, demanda nocional y efectiva de bienes; $L^{D,\mathrm{ef}}$, demanda efectiva de trabajo.}
\end{figure}

\subparagraph{c) Extensiones}

\begin{itemize}
\item El libro de \textcite{BarroGrossman1976Money} y los equilibrios con racionamiento de \textcite{Benassy1975Neo} y \textcite{Dreze1975Existence} generalizan el modelo; \autor{Malinvaud} lo convierte en una tipología del paro.
\end{itemize}

\subparagraph{d) Valoración}

\begin{itemize}
\item El propio \autor{Barro} lo abandonó: en «Second Thoughts on Keynesian Economics», \textcite{Barro1979Second} sostiene que unos precios fijos sin explicación dejan sin aprovechar ganancias del intercambio que los agentes aprovecharían, y se pasa a la nueva macroeconomía clásica \vertema{3.A.6}.
\end{itemize}

\paragraph{\autor{Malinvaud} (1977): los regímenes del paro}

\subparagraph{a) Idea}

\begin{itemize}
\item En Francia, los sindicatos decían que subir los salarios aumentaría el consumo y el empleo, y la patronal, que lo reduciría porque encarecería el trabajo. En \emph{The Theory of Unemployment Reconsidered}, \textcite{Malinvaud1977Theory} muestra que los 2 pueden tener razón: depende de qué mercado esté racionado.
\end{itemize}

\subparagraph{b) Modelo}

\begin{itemize}
\item \emph{i) Supuestos}: los de \autor{Barro} y \autor{Grossman}, con $W$ y $P$ fijos a corto plazo y una oferta de trabajo $\overline{L}$ dada.
\item \emph{ii) Desarrollo}: la producción es la menor de 3 cantidades, la demandada, la que las empresas quieren ofrecer a ese salario real y la de pleno empleo:
  \begin{equation}
  Y=\min\Big\{\,Y^{D}\Big(\overline{A},\frac{M}{P}\Big),\ \ Y^{S}\Big(\frac{W}{P}\Big),\ \ F(\overline{L})\,\Big\}, \label{eq:malinvaud}
  \end{equation}
  con $\overline{A}$ el gasto autónomo. Según cuál sea la menor, la economía está en uno de 3 \emph{regímenes} (gráfico~\ref{fig:malinvaud} y tabla~\ref{tab:malinvaud}):
  \begin{itemize}
  \item \concepto{paro keynesiano} ($Y=Y^D$): exceso de oferta de bienes y de trabajo; las empresas producirían más a ese salario real, pero no venden. Es el paro de \autor{Keynes};
  \item \concepto{paro clásico} ($Y=Y^S$): el salario real es tan alto que a las empresas no les compensa producir más, aunque podrían vender: exceso de demanda de bienes y exceso de oferta de trabajo;
  \item \concepto{inflación reprimida} ($Y=F(\overline{L})$): precios y salarios demasiado bajos para la cantidad de dinero; exceso de demanda en los 2 mercados.
  \end{itemize}
  Los 3 regímenes se tocan en el equilibrio walrasiano, $\mathcal{W}$, el único punto sin racionamiento.\footnote{Notación local: $\mathcal{W}$ caligráfica para el punto walrasiano, porque $W$ es el salario nominal; $Y^D$ e $Y^S$, la producción demandada y la ofrecida.}
\item \emph{iii) Implicaciones}: la política depende del diagnóstico:
  \begin{itemize}
  \item con paro keynesiano, más demanda, y subir los salarios puede crear empleo por el consumo: tienen razón los sindicatos;
  \item con paro clásico, más demanda no sirve: hay que bajar el salario real o subir la rentabilidad, y tiene razón la patronal;
  \item una misma tasa de paro puede pedir políticas opuestas;
  \item vista así, la revolución keynesiana es un cambio de diagnóstico: del paro clásico, el de \autor{Pigou}, al paro keynesiano, el de la Gran Depresión.
  \end{itemize}
\end{itemize}

\begin{figure}[H]
\centering
\caption{Los regímenes de \autor{Malinvaud} en el plano del salario nominal y el nivel de precios}\label{fig:malinvaud}
\grafico[0.55\textwidth]{malinvaud-regimenes}
\fuente{elaboración propia a partir de \textcite{Malinvaud1977Theory} y \textcite[295]{Tsoulfidis2010Competing}. K, paro keynesiano; C, paro clásico; I, inflación reprimida. La frontera entre K e I es horizontal: con el gasto autónomo, la cantidad de dinero y la oferta de trabajo dados, y si el consumo no depende del salario real, solo hay un nivel de precios con el que la demanda absorbe la producción de pleno empleo.}
\end{figure}

\begin{table}[H]
\centering
\caption{Los regímenes de \autor{Malinvaud} según el mercado racionado}\label{tab:malinvaud}
\small
\begin{tabular}{>{\raggedright\arraybackslash}m{3.2cm}>{\centering\arraybackslash}m{3cm}>{\centering\arraybackslash}m{3cm}>{\raggedright\arraybackslash}m{4.2cm}}
\toprule
\rowcolor{tceeazulpalido}\textbf{Régimen} & \textbf{Bienes} & \textbf{Trabajo} & \textbf{Política}\\
\midrule
\textcolor{tceeazul}{\textbf{Paro keynesiano}} & Exceso de oferta & Exceso de oferta & Más demanda\\
\addlinespace[0.3em]
\textcolor{tceeverde}{\textbf{Paro clásico}} & Exceso de demanda & Exceso de oferta & Salario real más bajo, más rentabilidad\\
\addlinespace[0.3em]
\textcolor{tceerojo}{\textbf{Inflación reprimida}} & Exceso de demanda & Exceso de demanda & Menos demanda\\
\addlinespace[0.3em]
\textcolor{tceedorado}{Subconsumo} & Exceso de oferta & Exceso de demanda & Caso aparte, fuera del modelo básico\\
\bottomrule
\end{tabular}
\fuente{elaboración propia a partir de \textcite{Malinvaud1977Theory} y \textcite{MuellbauerPortes1978Macroeconomic}.}
\end{table}

\subparagraph{c) Extensiones}

\begin{itemize}
\item \textcite{MuellbauerPortes1978Macroeconomic} añaden existencias y estudian un 4.º caso, el \emph{subconsumo} (exceso de oferta de bienes y de demanda de trabajo), que no es un régimen del modelo básico.
\item La tipología pasó al debate sobre el paro europeo y a la distinción entre paro keynesiano y clásico del tema de la tasa natural \vertema{3.A.28}.
\end{itemize}

\subparagraph{d) Evidencia empírica}

\begin{itemize}
\item En el proyecto que coordinaron \textcite{DrezeBean1990Europe} se estimaron modelos de desequilibrio para 10 países: en los 80, el paro europeo tenía una parte clásica, por los costes laborales tras las crisis del petróleo, y una parte keynesiana, por la demanda.
\end{itemize}

\subparagraph{e) Valoración}

\begin{itemize}
\item \emph{Virtudes}: una tipología clara, que todavía sirve para diagnosticar el paro y elegir la política.
\item \emph{Críticas}: los precios y los salarios son exógenos y el modelo es estático: no dice cuánto dura cada régimen ni cómo se pasa de uno a otro.
\end{itemize}

\paragraph{Valoración y legado de los keynesianos del desequilibrio}

\begin{itemize}
\item \emph{Virtudes}: dieron a \autor{Keynes} fundamentos microeconómicos; hicieron del paro un problema de coordinación, la idea que retoman los modelos con equilibrios múltiples \parencite{CooperJohn1988Coordinating} \vertema{3.A.21}; y distinguieron tipos de paro que piden políticas distintas.
\item \emph{Críticas}: no explican por qué los precios no se ajustan \parencite{Barro1979Second}; sus expectativas no son racionales y sus parámetros cambiarían con la política, la crítica de \autor{Lucas} \vertema{3.A.6}; y los modelos con racionamiento se complican mucho fuera de los casos sencillos.
\item \emph{Lo que añaden, en 2 gráficos}: con \autor{Patinkin}, el paro involuntario es un tránsito y la economía vuelve sola al pleno empleo (gráfico~\ref{fig:patinkin}); con \autor{Clower} y \autor{Barro} y \autor{Grossman}, es un equilibrio con racionamiento del que solo la saca la demanda (gráfico~\ref{fig:racionamiento}).
\item \emph{Legado}: la nueva economía keynesiana conserva su pregunta y responde a la crítica: explica con microfundamentos por qué los precios y los salarios son rígidos (costes de menú, salarios de eficiencia, contratos) \vertema{3.A.7}.
\end{itemize}

\begin{table}[H]
\centering
\caption{\autor{Keynes} y sus herederos: por qué hay paro y qué hacer}\label{tab:herederos}
\small
\begin{tabular}{>{\raggedright\arraybackslash}m{2.3cm}>{\raggedright\arraybackslash}m{2.75cm}>{\raggedright\arraybackslash}m{2.75cm}>{\raggedright\arraybackslash}m{2.75cm}>{\raggedright\arraybackslash}m{2.75cm}}
\toprule
\rowcolor{tceeazulpalido} & \textbf{\autor{Keynes}} & \textbf{Síntesis neoclásica} & \textbf{Postkeynesianos} & \textbf{Desequilibrio}\\
\midrule
¿Por qué hay paro? & Demanda efectiva insuficiente & Salario nominal rígido & Demanda efectiva, también a largo plazo & Precios que no vacían: fallo de coordinación\\
\addlinespace[0.3em]
Precios y salarios & Salario nominal rígido; bajarlo no cura & Rígidos a corto, flexibles a largo & Margen sobre costes; reparto según el poder & Fijos a corto plazo, sin explicar\\
\addlinespace[0.3em]
Dinero & Exógeno; preferencia por la liquidez & Exógeno; neutral a largo plazo & Endógeno; fragilidad financiera & Saldos reales y racionamiento\\
\addlinespace[0.3em]
Incertidumbre & Central & Secundaria & Central & Secundaria\\
\addlinespace[0.3em]
Política & Fiscal; socialización de la inversión & Sintonía fina \vertema{3.A.5} & Fiscal, de rentas y regulación financiera & Según el régimen\\
\bottomrule
\end{tabular}
\fuente{elaboración propia.}
\end{table}

\begin{itemize}
\item \emph{Recapitulación}: los postkeynesianos rechazan el instrumental neoclásico para quedarse con la demanda efectiva y la incertidumbre; los keynesianos del desequilibrio lo aceptan para encontrar en él el paro involuntario. Frente a la síntesis, los 2 coinciden en que el paro no se explica solo con un salario que no baja.
\end{itemize}

% ==========================================================================
\section*{Conclusión}
\minutos{1}

\begin{itemize}
\item \textbf{\emph{Recapitulación}}:
  \begin{itemize}
  \item En conclusión, \autor{Keynes} invierte la causalidad del modelo neoclásico: la demanda efectiva fija la producción y el empleo, el tipo de interés es el precio de la liquidez y no el del ahorro, y bajar los salarios no garantiza el pleno empleo.

  \item Sus herederos no se quedaron con la lectura de la síntesis: los postkeynesianos de Cambridge llevaron sus ideas a la distribución, al capital y a las finanzas, y los teóricos del desequilibrio explicaron el paro como un fallo de coordinación entre mercados que no se vacían.
  \end{itemize}
\item \textbf{\emph{Relevancia}}:
  \begin{itemize}
  \item La Gran Depresión le dio la razón a \autor{Keynes}; la estanflación de los 70, a sus críticos; la Gran Recesión y la pandemia devolvieron a primer plano los multiplicadores, la trampa de la liquidez y el «momento \autor{Minsky}».
  \end{itemize}
\item \textbf{\emph{Extensiones}}:
  \begin{itemize}
  \item La síntesis y su crítica monetarista \vertema{3.A.5}; la nueva macroeconomía clásica \vertema{3.A.6} y la nueva economía keynesiana, que fundamenta las rigideces que el desequilibrio daba por supuestas \vertema{3.A.7}; la demanda de dinero \vertema{3.A.35}; el ciclo financiero \vertema{3.A.42}; el crecimiento y la distribución \vertema{3.A.43} \vertema{3.A.45}.
  \end{itemize}
\item \textbf{\emph{Opinión}}:
  \begin{itemize}
  \item El gran acierto de \autor{Keynes} fue la pregunta, si el paro puede ser un fallo del sistema y no de un mercado.

  \item La macroeconomía de hoy la responde con modelos que él no habría escrito, pero la respuesta práctica de 2020 fue la suya \parencite{CE2026NGEU}.
  \end{itemize}
\item \textbf{\emph{Idea final}}:
  \begin{itemize}
  \item A la pregunta de si una economía con paro se corrige sola, \autor{Keynes} responde que quizá sí, pero a largo plazo, y que mientras tanto el paro de hoy es hambre de hoy.
  \item Sus herederos añaden que el largo plazo se construye con la demanda de hoy: la inversión que no se hace hoy es la capacidad que faltará mañana.
  \item Por eso el pan de hoy no tiene por qué ser hambre para mañana: si se invierte, es lo que permite construir el futuro sin el estómago vacío.
  \end{itemize}
\end{itemize}

\clearpage
\appendix

\section{Anexo: los modelos neoclásico y keynesiano, mercado a mercado}\label{anexo:modelos}

\begin{itemize}
\item Este anexo desarrolla lo que el apartado~\ref{sec:neoclasicos} y el bloque~\ref{sec:formalizacion} resumen: el punto de partida marshalliano y, mercado a mercado, las piezas de los 2 modelos.
\end{itemize}

\subsection{El mercado de pescado de \autor{Marshall}}\label{anexo:marshall}

\begin{itemize}
\item \autor{Keynes} se formó con \autor{Marshall} y razona como él: equilibrio parcial, periodos definidos por el tiempo que tienen los agentes para ajustarse y expectativas a corto plazo \parencite[libro~V, cap.~V]{Marshall1890Principles} \vertema{3.A.16}.
\item En su mercado de pescado (gráfico~\ref{fig:marshall-periodos}) hay 2 ofertas:
  \begin{itemize}
  \item la \emph{oferta de mercado} de cada día ($\mathrm{MS}$) es vertical: el pescado ya está en la lonja;
  \item la \emph{oferta normal} ($\mathrm{NS}$) es creciente: con tiempo, entran más barcos y la cantidad se ajusta.
  \end{itemize}
\item Si sube la demanda, el primer día el precio salta hasta vaciar la lonja (equilibrio temporal, $B$), con ganancias extraordinarias para los que ya pescaban; con el tiempo llegan más barcos y el precio baja hacia el equilibrio normal ($C$). En ningún momento queda pescado sin vender: los precios se ajustan antes que las cantidades.
\item \autor{Keynes} conserva el corto plazo y las expectativas, pero cambia el orden: en su economía, a corto plazo se ajustan antes las cantidades, y un mercado puede no vaciarse. Es la lectura que hará \autor{Leijonhufvud} (\ref{sec:desequilibrio}).
\end{itemize}

\begin{figure}[H]
\centering
\caption{El mercado de pescado de \autor{Marshall}: equilibrio temporal ($B$) y equilibrio normal ($C$) tras un aumento de la demanda}\label{fig:marshall-periodos}
\grafico[0.5\textwidth]{marshall-periodos}
\fuente{elaboración propia a partir de \textcite[fig.~1.1]{DeVroey2016History} y \textcite[libro~V, cap.~V]{Marshall1890Principles}.}
\end{figure}

\subsection{El mercado de trabajo}\label{anexo:trabajo}

\begin{itemize}
\item \emph{Neoclásicos}: una empresa representativa elige el trabajo para maximizar su beneficio con el capital dado,
  $$\optim{\max}{L}{P\,F(L,\overline{K})-W\,L}{L\geq 0}$$
  y la condición de primer orden, $\mathrm{PMg}_L=W/P$, da una demanda de trabajo decreciente en el salario real. Los hogares eligen entre consumo y ocio, y si domina el efecto sustitución ofrecen más trabajo cuanto mayor es el salario real \vertema{3.A.25}. El salario real se ajusta hasta vaciar el mercado; los cambios de precios se trasladan en la misma proporción al salario nominal y el empleo no cambia: la oferta agregada es vertical.
\item \emph{\autor{Keynes}}: la misma demanda de trabajo, pero con el salario nominal negociado, $W_0$, y el empleo fijado por la demanda de bienes. Si los precios suben, el salario real baja y se contrata más (oferta agregada creciente mientras haya recursos ociosos); si bajan, se contrata menos y aparece el paro involuntario (gráfico~\ref{fig:keynesiano-shock}).
\end{itemize}

\subsection{El mercado de bienes}\label{anexo:bienes}

\begin{itemize}
\item \emph{Neoclásicos}: de $Y=C+I+G$ y del ahorro privado, $S_{\mathrm{priv}}=Y-T-C$, sale
  \begin{equation}
  S_{\mathrm{priv}}+(T-G)=I: \label{eq:anexo-ahorro}
  \end{equation}
  el ahorro privado más el público financia la inversión. Con el presupuesto equilibrado, el ahorro privado iguala la inversión, y el tipo de interés real los iguala (gráfico~\ref{fig:fondos-prestables}). Si el Estado gasta más, sube el tipo y el gasto privado baja en la misma cuantía.
\item \emph{\autor{Keynes}}: el multiplicador sale de sumar las rondas de gasto. Un aumento del gasto autónomo, $\Delta\overline{A}$, es renta de alguien, que gasta la parte $c\,(1-t)$; esa parte es renta de otro, que gasta otra vez esa proporción, y así sucesivamente:
  \begin{equation}
  \Delta Y=\Delta\overline{A}\,\Big[1+c\,(1-t)+\big(c\,(1-t)\big)^2+\cdots\Big]=\frac{\Delta\overline{A}}{1-c\,(1-t)}, \label{eq:anexo-multiplicador}
  \end{equation}
  una serie geométrica que converge porque $0<c\,(1-t)<1$. Es el mismo resultado que \eqref{eq:multiplicador}. \findemo
\end{itemize}

\subsection{El mercado de dinero}\label{anexo:dinero}

\begin{itemize}
\item \emph{Neoclásicos: la teoría cuantitativa} \vertema{3.A.35}:
  \begin{itemize}
  \item \autor{Fisher}: la ecuación de cambio, $M\,V\equiv P\,T$, con $T$ el volumen de transacciones (no los impuestos), es una identidad; se convierte en teoría con 3 supuestos: $V$ depende de hábitos de pago e instituciones que cambian despacio, $T$ (o $Y$) lo fija el sector real y $M$ es exógena. Entonces la causalidad va de $M$ a $P$ \parencite{Fisher1911Purchasing};
  \item Cambridge: la demanda de saldos de caja es una fracción $k$ de la renta nominal, $M^D=k\,P\,Y$; \autor{Marshall} y \textcite{Pigou1917Value} la ven como una decisión de los individuos, que podría depender del tipo de interés y de las expectativas, aunque la tratan como constante. Con $k=1/V$, las 2 ecuaciones coinciden.\footnote{La idea de que más dinero rebaja su valor es muy anterior: \autor{Martín de Azpilcueta} la expuso en 1556, en la Escuela de Salamanca, 12 años antes que \autor{Jean Bodin} \parencites{Azpilcueta1556Cambios}{Bodin1568Reponse}.}
  \end{itemize}
\item \emph{\autor{Keynes}}: la demanda de Cambridge con un término más, el especulativo, que depende del tipo de interés: $M/P=k\,Y-h\,i$ \eqref{eq:liquidez}. Si $h=0$ vuelve la de Cambridge; cuanto mayor es $h$, menos determina el dinero los precios y más el tipo de interés.
\end{itemize}

\subsection{Las políticas en los 2 modelos}\label{anexo:politicas}

\begin{itemize}
\item \emph{Expansión fiscal}: en el modelo neoclásico, la demanda de fondos se desplaza, sube el tipo de interés real y la producción no cambia; en el de \autor{Keynes}, con recursos ociosos, la producción sube con el multiplicador \eqref{eq:anexo-multiplicador}.
\item \emph{Expansión monetaria}: en el neoclásico, la demanda agregada cuantitativa se desplaza sobre la oferta vertical y solo suben los precios; en el de \autor{Keynes}, baja el tipo de interés, sube la inversión y la producción, salvo con trampa de la liquidez o inversión insensible.
\item \emph{Rebaja salarial}: en el neoclásico, devuelve el pleno empleo; en el de \autor{Keynes}, solo crea empleo si sube la inversión, y nada lo garantiza (capítulo~19). El resumen es la tabla~\ref{tab:politicas}.
\end{itemize}

\section{Anexo: «clásicos» y «neoclásicos» en la \emph{Teoría general}}\label{anexo:clasicos}

\begin{itemize}
\item \autor{Keynes} da a «clásicos» un sentido propio. En una nota del capítulo~1 explica que el nombre lo inventó \autor{Marx} para \autor{Ricardo}, \autor{James Mill} y sus predecesores, y que él lo extiende a quienes adoptaron y perfeccionaron la teoría ricardiana, como \autor{John Stuart Mill}, \autor{Marshall}, \autor{Edgeworth} y \autor{Pigou} \parencite[cap.~1]{Keynes1936}. Hoy llamamos neoclásicos a casi todos estos últimos \vertema{3.A.2} \vertema{3.A.3}.
\item Lo que los une, para \autor{Keynes}, no es la teoría del valor, sino la ley de \autor{Say} y la confianza en que, con salarios flexibles, la economía tiende al pleno empleo.
\item La agrupación es discutible: \autor{Ricardo} aceptaba la ley de \autor{Say}, pero en el capítulo sobre la maquinaria de la 3.ª edición de sus \emph{Principios} admitió que las máquinas pueden dejar sin trabajo a los obreros; no daba por supuesto el pleno empleo \parencite{Astarita2009Keynes}.
\item El término «neoclásico» lo acuñó \autor{Thorstein Veblen} hacia 1900 para subrayar la continuidad de \autor{Marshall} con los clásicos \parencite{Aspromourgos1986Neoclassical}.
\end{itemize}

\section{Anexo: notas del tema anterior}\label{anexo:notas}

\begin{itemize}
\item Este anexo conserva notas del tema anterior que no se cantan, pero amplían lo que el tema cuenta.
\end{itemize}

\begin{anotaciones}
\begin{itemize}
\item \emph{Neutralidad, superneutralidad y dicotomía}:
  \begin{itemize}
  \item el dinero es \emph{neutral} si los cambios en su cantidad no tienen efectos reales, y \emph{superneutral} si tampoco los tienen los cambios en su tasa de crecimiento (en un modelo estático, como el del apartado~\ref{sec:neoclasicos}, no tiene sentido);
  \item hay \emph{dicotomía clásica} si las variables reales (empleo y producción) se determinan en un subsistema independiente de la cantidad de dinero y del nivel de precios: el dinero es un velo;
  \item son conceptos distintos: \textcite{Sargent1979Macroeconomic} construye un ejemplo en el que se cumple la dicotomía y no la neutralidad, aunque en la práctica conducen a definiciones parecidas.
  \end{itemize}
\item \emph{\autor{David Hume} y el corto plazo}: \textcite{Hume1752Money} distinguía el largo plazo, en el que más dinero sube los precios en la misma proporción, del corto, en el que, mientras los precios se ajustan, el dinero nuevo anima la industria y el comercio. Los clásicos buscaban leyes generales y se quedaron con el largo plazo; \autor{Keynes} vuelve al corto.
\item \emph{La causalidad en el mercado de trabajo}: en el modelo neoclásico, la causalidad va del salario a la demanda de trabajo, que se fija sin mirar la demanda de los productos de la empresa; \autor{Keynes} la invierte: la demanda de bienes fija el empleo, y el empleo, el salario real \parencite{Gali2013Keynes}.
\item \emph{Otros escritos}: los ensayos de \emph{Essays in Persuasion} \parencite{Keynes1931Essays} y las biografías de economistas y políticos de \emph{Essays in Biography} \parencite{Keynes1933Essays}.
\item \emph{Los Nobel de los seguidores de la síntesis} \parencite{Nobel2026Economia}:
  \begin{itemize}
  \item \autor{Jan Tinbergen}, con \autor{Ragnar Frisch} (1969), por los modelos dinámicos para analizar los procesos económicos;
  \item \autor{Paul Samuelson} (1970), por elevar el nivel del análisis en la ciencia económica;
  \item \autor{John Hicks}, con \autor{Kenneth Arrow} (1972), por la teoría del equilibrio general y del bienestar;
  \item \autor{James Meade}, con \autor{Bertil Ohlin} (1977), por la teoría del comercio internacional;
  \item \autor{Lawrence Klein} (1980), por los modelos econométricos y su aplicación a las fluctuaciones y la política económica;
  \item \autor{Franco Modigliani} (1985), por el análisis del ahorro y de los mercados financieros;
  \item \autor{Robert Solow} (1987), por la teoría del crecimiento económico.
  \end{itemize}
\item \emph{El nombre de la síntesis}: \autor{Paul Samuelson} llamó «síntesis neoclásica» a la combinación de lo que creía valioso de la economía anterior y de la teoría de la renta keynesiana, en la 3.ª edición de su manual (1955). Para los keynesianos de los años 40 resultó demasiado neoclásica, y para los monetaristas de los 60, demasiado keynesiana \vertema{3.A.5}.
\item \emph{Leer hoy la \emph{Teoría general}}: para \textcite{Mankiw2006Macroeconomist}, la experiencia es a la vez estimulante y frustrante: es la obra de una gran mente ante un problema social enorme, pero como argumento lógico parece incompleta.
\end{itemize}
\end{anotaciones}

\section{Anexo: aspectos para el test}\label{anexo:test}

\begin{anotaciones}
Preguntas tipo test del primer ejercicio: \href{https://www.victorgutierrezmarcos.es/oposicion/temario/primer-ejercicio/test/simulador.html}{simulador de test}.
\end{anotaciones}

\bibliografia
\end{document}
